\documentclass[openany]{book}

\usepackage{cianatex}

\usepackage{cianacolors}

\usepackage[cianabook]{cianatheorems}

\addbibresource{MeccanicaAnaliticaSchenkel.bib}

\newcommand\GL{\operatorname{GL}}
\newcommand\affcat{\catname{Aff}}
\renewcommand\C{\mc{C}}
\renewcommand\A{\mc{A}}
\renewcommand\V{\mathbb{V}}

\title{Fondamenti di Fisica Matematica 1}
\author{F. Troncana, dalle note del prof. A. Schenkel}
\date{A.A. 2025/2026}

\begin{document}

\maketitle

\tableofcontents

\chapter{Geometria}

In questo capitolo daremo un'esposizione relativamente minimale degli strumenti di geometria che ci serviranno per trattare la meccanica classica in modo soddisfacente ai fini del nostro corso.

\section{Aspetti fondamentali della Geometria}

Esistono diversi tipi di spazi che svolgono un ruolo importante nella meccanica classica, ad esempio:\begin{itemize}
    \item Spazi vettoriali;
    \item Spazi affini;
    \item Spazi topologici;
    \item Varietà
\end{itemize}
In questa prima sezione ci ricorderemo alcuni fatti su spazi con cui il lettore dovrebbe essere già familiare e ne introdurremo di nuovi che useremo in seguito.
\paragraph{Nota bene:} tutti gli spazi considerati saranno definiti sul campo $\R$ dei numeri reali ai fini delle applicazioni alla meccanica classica. 

\subsection{Spazi vettoriali}

Assumiamo che il lettore abbia già familiarità con il concetto \bemph{($\R$-)spazio vettoriale}\footnote{Ovvero, un gruppo abeliano dotato di un'azione canonica di $\R$ compatibile con la struttura di gruppo.} e di morfismi tra spazi vettoriali, ovvero le \bemph{mappe lineari}. Ricordiamo alcuni fatti che useremo più avanti:

\begin{example}{Lo spazio di dimensione finita $\R^n$}{}
    L'esempio standard di spazio vettoriale è $\R^n$, con $n \in \N$.\\
    Per ogni $n, m \in \N$ esiste\footnote{una volta scelta una base} una biezione $\M_{n\times m}(\R) \iso \hom(\R^m,\R^n)$ che associa una matrice $M$ alla mappa lineare $f_M : \bar{v}\mapsto M\bar{v}$
\end{example}

Notiamo che se $V$ è uno spazio vettoriale di dimensione $n \in \N$, esiste un\footnote{non unico!} isomorfismo $\R^n \iso V$.\\
Infatti scegliendo una base $\{e_i\in V\}_{1\le i \le n}$ un isomorfismo lineare può essere definito da 
\[ e: \bar{v} = \begin{pmatrix} v^1 \\ \vdots \\ v^n \end{pmatrix}\mapsto \sum_{i=1}^n v^i e_i  \]
In particolare, dato che per due basi $\{e_i\}_i$, $\{e'_i\}_i$ esiste una matrice $B = (B_i^j)$ in $\GL(n)$ tale che 
\[ e_i = \sum_{j = 1}^n B^j_i e'_i \]
Otteniamo un automorfismo lineare
\[\begin{tikzcd}
	& V & \\
	{\R^n} && {\R^n}
	\arrow["e"{description}, from=2-1, to=1-2]
	\arrow["B"{description}, from=2-1, to=2-3]
	\arrow["{e'}"{description}, from=2-3, to=1-2]
\end{tikzcd}\]
Nel contesto della fisica interpretiamo $e$ come una \bemph{scelta di coordinate lineari} per $V$ e $B$ come un \bemph{cambiamento di coordinate}.\\
Potremmo essere tentati di pensare che lo spaziotempo della fisica classica sia descritto da spazi vettoriali $V \cong \R^{3+1}$, con tre dimensioni spaziali e una temporale, ma non è così in quanto ogni spazio vettoriale ha un punto privilegiato, l'origine $0$: "dimenticare" questa origine ci porta al concetto di spazio affine.

\subsection{Spazi affini}

\begin{definition}{Spazio affine}{spazio affine}
    Si dice \bemph{spazio affine} su uno spazio vettoriale $V$ una coppia $(A,+)$, dove $A$ è un insieme e $+:A\times V\to A$ una mappa tale che:\begin{enumerate}
        \item $+$ è un'azione destra di $V$ su $A$, ovvero tale che per ogni $P\in A$ e $v,w\in V$ valgano \[ (P+v)+w = P+(v+w),\qquad P+0=P \]
        \item L'azione è libera, ovvero $P+v = P$ se e solo se $v=0$.
        \item L'azione è transitiva, ovvero per ogni $P,Q \in A$ esiste un vettore $v \in V$ tale che $P = Q+v$, lo indicheremo con $v = P-Q$.
    \end{enumerate}
    Un morfismo di spazi affini $A,A'$ è una mappa $\psi: A\to A'$ tale che esista una mappa lineare $\di\psi : V\to V'$ tale che per ogni $P \in A$ e $v \in V$ valga 
    \[ \psi(P+v) = \psi(P)+\di\psi(v) \]
\end{definition}

È facile dimostrare che esiste una mappa $- : A\times A \to V$ che manda $(P,Q)$ in $v = P-Q$. È possibile dare una definizione equivalente di spazio affine in termini degli assiomi che deve rispettare questa mappa, come in \cite{Moretti2024}.\\
Notiamo anche che la mappa lineare $\di\psi$ è determinata univocamente dalla mappa $\psi$: infatti vediamo che se $\delta\psi$ fosse un'altra mappa lineare che soddisfi gli assiomi, varrebbe
\[ \psi(P+v) = \psi(P)+ \delta\psi(v) = \left(\psi(P)+ \di\psi(v)\right) + \left(\delta\psi(v) - \di\psi(v)\right) \]
Che implica $\delta\psi(v) \equiv \di\psi(v)$.

\begin{example}{Struttura affine dello spazio vettoriale}{}
    Notiamo che munendo $V$ dell'azione canonica destra su sè stesso data dalla somma di vettori $(v)+w = (v+w)$, questo assume una struttura di spazio affine: in particolare ogni ogni mappa lineare $L:V\to V'$ definisce una mappa affine tale che $\di L = L$.\\
    Non tutte le mappe affini $\psi:V\to V'$ tuttavia hanno questa forma: per definizione abbiamo
    \[ \psi(v) = \psi(0+v) = \psi(0) + \di\psi(v) \]
    Dunque notiamo che abbiamo una biezione $\hom_\affcat(V,V') \iso V'\times\hom_\veccat(V,V') $ che manda una mappa affine $\psi$ nella coppia $(\psi(0),\di\psi)$.\\
    In parole povere, una mappa affine è data da una traslazione e una mappa lineare: la mappa $\psi$ è lineare se e solo se $\psi(0) = 0'$.
\end{example}

\begin{proposition}{Isomorfismo con lo spazio sottostante}{}
    Sia $A$ uno spazio affine su $V$ e muniamo $V$ della sua struttura canonica di spazio affine: allora esiste un\footnote{non unico} isomorfismo di spazi affini $V\iso A$.
    \proof 
    Scegliamo un punto $P_0 \in A$ e definiamo la mappa $\psi: v\mapsto P_0+v$ e notiamo che è invertibile con inversa data da $\psi^{-1}(P) = P-P_0$.\\
    La mappa $\psi$ è evidentemente affine, come lo è $\psi^{-1}$, segue la tesi.
    \qed 
\end{proposition}

Come conseguenza diretta della proposizione precedente, ogni spazio affine $A$ di dimensione $n\in \N$ è dotato di un\footnote{non unico} isomorfismo affine con $\R^n$: scegliendo un punto base $P_0 \in A$ e una base $\{e_i\}_i$ di $V$ questo è dato da
\[ (P_0,e): \bar{v} = \begin{pmatrix} v^1\\\vdots\\v^n \end{pmatrix}\mapsto P_0 + \sum_{i=1}^n v^ie_i\]
In particolare, considerando due punti $P_0,P_0' \in A$ e dato che per due basi $\{e_i\}_i$, $\{e'_i\}_i$ esiste una matrice $B = (B_i^j)$ in $\GL(n)$ tale che 
\[ e_i = \sum_{j = 1}^n B^j_i e'_i \]
Otteniamo un automorfismo affine 
\[\begin{tikzcd}
	& A & \\
	{\R^n} && {\R^n}
	\arrow["{(P_0,e)}"{description}, from=2-1, to=1-2]
	\arrow["{(P_0'-P_0, B)}"{description}, from=2-1, to=2-3]
	\arrow["{(P'_0,e')}"{description}, from=2-3, to=1-2]
\end{tikzcd}\]
Gli spazi affini possono essere modelli utili per alcuni fenomeni fisici.\\
Modellando il tempo come uno spazio $A^1$ di dimensione $1$ e lo spazio come uno spazio $A^3$ di dimensione $3$, le mappe affini $\psi: A^1\to A^3$ descrivono il moto uniforme di una particella su rette affini: infatti scegliendo qualsiasi coordinate, queste mappe assumono la forma esplicita
\[\begin{tikzcd}
	{A^1} && {A^3} \\
	{\R^1} && {\R^3} \\
	t && {\bar{x} + \bar{v}t}
	\arrow["\psi"{description}, from=1-1, to=1-3]
	\arrow[from=2-1, to=1-1]
	\arrow[from=2-1, to=2-3]
	\arrow[from=2-3, to=1-3]
	\arrow[maps to, from=3-1, to=3-3]
\end{tikzcd}\]
Dove $\bar{v}$ è la velocità dela particella e $\bar{x}_0$ la sua posizione al tempo $t=0$.\\
Tuttavia, per descrivere moti più interessanti (ad esempio quando le particelle sono sottoposte a delle forze) è necessario passare ad un formalismo ancora più generale.

\subsection{Varietà lisce}

Per ottenere un formalismo geometrico più potente e flessibile, introduciamo il concetto di \bemph{varietà liscia}. Rinfreschiamoci la memoria su alcuni concetti preliminari:

\begin{definition}{Spazio topologico}{}
    Si dice \bemph{spazio topologico} un insieme $X$ dotato di un \bemph{frame} di suoi sottoinsiemi $\tau\subset 2^X$, detto \bemph{topologia} su $X$, ovvero una famiglia di suoi sottoinsiemi detti \bemph{aperti} tale che:\begin{enumerate}
        \item $\varnothing \in \tau$ e $X\in \tau$;
        \item Data una qualsiasi sottofamiglia $\{U_i\}_i$ di $\tau$, la sua unione $\bigcup_i U_i$ appartiene a $\tau$;
        \item Data una sottofamiglia \textit{finita} $\{U_i\}_i$ di $\tau$, la sua intersezione $\bigcap_i U_i$ appartiene a $\tau$, o equivalentemente dati due aperti $U$ e $V$ di $\tau$ la loro intersezione appartiene a $\tau$.
    \end{enumerate}
    Un morfismo di spazi topologici $f: (X,\tau)\to(Y,\xi)$ è dato da una mappa di insiemi $f:X\to Y$ tale che la sua immagine inversa $f^{-1}: 2^Y \to 2^X$ si restringa ad una mappa di frame $\xi\to\tau$, ovvero tale che la controimmagine attraverso $f$ di un aperto $U \in \xi$ sia un aperto di $\tau$.\\
    Un isomorfismo di spazi topologici si dice omeomorfismo.
\end{definition}

Indicheremo la topologia euclidea su $\R^n$ con $\tau_e^n$ e assumeremo che $\R^n$ e i suoi sottospazi ne siano sempre dotati.

\begin{definition}{Funzione liscia nel caso reale}{}
    Una funzione $f: U\to W$, dove $U\in \tau_e^n$ e $W \in \tau_e^m$ con $n,m \in \N$ si dice \bemph{liscia} o $\C^\infty$ se è differenziabile infinite volte.
\end{definition}

Assumiamo che il lettore abbia familiarità con alcuni concetti di topologia generale dal primo modulo del corso Geometria B

\begin{definition}{Varietà topologica, differenziabile e liscia}{varietà}
    Sia $(M,\mu)$ uno spazio topologico di Hausdorff\footnote{Ogni coppia di punti distinti è contenuta in due aperti disgiunti} a base numerabile\footnote{Esiste una sottofamiglia numerabile $B$ di aperti tale che ogni aperto di $\mu$ sia unione di aperti di $B$}. \begin{enumerate}
        \item Dato un aperto $U \in \mu$, una \bemph{carta locale $n$-dimensionale} su $U$ è una mappa $\phi: U \to \R^n$ tale che $\phi(U)$ sia un aperto di $\R^n$ e la restrizione $\psi: U\iso\phi(U)$ sia un omeomorfismo.
        \item Due carte locali $n$-dimensionali $(U,\phi)$ e $(V,\psi)$ si dicono $\C^k$-compatibili per $k \in \N\cup\{\infty\}$ se le \bemph{mappe di trasferimento} $\psi \circ \phi^{-1}: \phi(U\cap V)\to\psi(U\cap V)$ e $\phi\circ\psi^{-1}:\psi(U\cap V)\to\phi(U\cap V)$ sono entrambe $\C^k$.
        \item Un \bemph{atlante di classe $\C^k$ di dimensione $n$} su $M$ è una famiglia di carte $n$-dimensionali $\{(U_i,\psi_i)\}_i$ tale che $\{U_i\}_i$ sia un ricoprimento di $M$ e le carte siano tutte $\C^k$-compatibili a due a due; questo si dice \bemph{massimale} se contiene ogni carta compatibile a ogni sua carta.
    \end{enumerate}
    $M$ si dice \bemph{varietà topologica, differenziabile di classe $\C^k$ o liscia di dimensione $n$} se possiede un atlante massimale rispettivamente di classe $\C^0$, $\C^k$ con $k \in \N_{>0}$ o $\C^\infty$ di dimensione $n$.\\
    Un morfismo di varietà di classe $\C^k$ $f:(M,\A)\to (M',\A')$ di dimensioni $n,n' \in\N$ è dato da una mappa continua $f:M\to N$ tale che per ogni $(U,\phi)\in \A$ e $(U',\phi') \in \A'$ tale che $f(U)\subset U'$, la composizione definita dal diagramma sia $\C^k$
    \[\begin{tikzcd}
    	& {M\ni U} && {U'\in M'} & \\
	    \\
	    {\R^n\ni \phi(U)} &&&& {\phi'(U')\in\R^{n'}}
	    \arrow["f"{description}, from=1-2, to=1-4]
	    \arrow["\phi"{description}, from=1-2, to=3-1]
	    \arrow["{\phi'}"{description}, from=1-4, to=3-5]
	    \arrow["{\phi'\circ f\circ \phi^{-1}}"{description}, from=3-1, to=3-5]
    \end{tikzcd}\]
\end{definition}
%Qui sarebbe bello aggiungere l'immagine di una varietà ma non ho voglia di farla in TIKZ
Una carta locale $(U,\phi)$ definisce su $M$ delle \bemph{coordinate locali} date dall'immagine di un punto in $U$ attraverso $\phi$
\[ \phi(P) = \begin{pmatrix} \phi^1(P)\\ \vdots \\ \phi^n(P) \end{pmatrix} \]
Queste ci permettono di definire un concetto di differenziazione locale su $M$\footnote{limitato dalla classe di differenziabilità di $M$ ovviamente}, invariante per transizione a qualsiasi altra carta compatibile.\\
Un atlante massimale di dimensione $n$ per $M$ in generale contiene una quantità mostruosa di carte, comparabile alla cardinalità di $\R^n$, ma nella pratica è sufficiente esibirne molte di meno per descrivere un unico atlante massimale, come vediamo nel prossimo risultato.

\begin{proposition}{Unicità dell'atlante massimale}{}
    Sia $\A$ un atlante su $M$ di dimensione $n \in \N$ di classe $\C^k$ con $k \in \N\cup\{\infty\}$. Questo è contenuto in un unico atlante massimale $\bar\A$, detto \bemph{atlante massimale indotto} da $\A$.
    \proof 
    Definiamo $\bar\A$ come l'insieme delle carte locali su $\M$ compatibili con tutte le carte locali di $\A$: chiaramente $\A\subset\bar\A$ e contiene per definizione tutte le carte compatibili a quelle di $\A$; l'unica cosa da dimostrare è che le carte di $\bar\A$ sono compatibili tra loro\footnote{ovvero che la compatibilità è una relazione transitiva}.\\
    Siano $(U,\phi)$ e $(V,\psi)$ due carte di $\bar\A$. Se $U\cap V=\varnothing$ non abbiamo nulla da dimostrare, dunque assumiamo che non sia vuoto. Per ogni $p \in U\cap V$ esiste dunque una qualche carta $(W,\chi)$ di $\A$ che lo contiene. La restrizione 
    \[\begin{tikzcd}
    	{\phi(U\cap V\cap W)} && {\psi(U\cap V \cap W)} \\
	    \\
	    & {\chi(U\cap V\cap W)}
	    \arrow["{\psi\circ\phi^{-1}}"{description}, from=1-1, to=1-3]
	    \arrow["{\chi\circ\phi^{-1}}"{description}, from=1-1, to=3-2]
	    \arrow["{\psi\circ\chi^{-1}}"{description}, from=3-2, to=1-3]
    \end{tikzcd}\]
    è differenziabile per compatibilità con $(W,\chi)$, dunque procedendo per esaustione sui punti di $U\cap V$ otteniamo la compatibilità delle due carte.\\
    L'unicità segue dalla massimalità.
    \qed
\end{proposition}

Vediamo che è sufficiente controllare la differenziabilità su un atlante non massimale:

\begin{proposition}{Incollamento per mappe differenziabili}{incollamento per mappe differenziabili}
    Sia $f:M\to M'$ una mappa continua e siano $\B,\B'$ due atlanti di classe $\C^k$ e dimensioni $n,n'$, non necessariamente massimali, che inducano le strutture differenziabili $\A := \bar\B$ e $\A':=\bar\B'$.\\
    Se la mappa definita dal diagramma della definizione \ref{def:varietà} è $\C^k$ per ogni coppia di carte in $\B,\B'$ che ne soddisfino le ipotesi, allora lo è anche per ogni coppia di carte in $\A,\A'$ che ne soddisfino le ipotesi, ovvero $f$ è un morfismo di varietà $\C^k$.
    \proof 
    Presa una coppia di carte $(U,\phi)$ e $(U',\phi')$ con $f(U)\subset U'$, la dimostrazione segue per esaustione sui punti di $U$ restringendo a intersezioni con carte $(V,\psi)$ e $(V',\psi')$ di $\B$ e $\B'$ analogamente alla proposizione precedente.
    \qed
\end{proposition}

Vediamo che sia gli spazi vettoriali che gli spazi affini sono esempi di varietà lisce:

\begin{example}{Spazio vettoriale come varietà liscia}{}
    Come visto sopra, posto uno spazio vettoriale $V$ di dimensione $n$ la scelta di una base $\{e_i\}_i$ induce un isomorfismo $e:\R^n\iso V$: possiamo usarlo per trasferire la topologia euclidea di $\R^n$ su $V$, ovvero l'unica topologia su $V$ tale che $e$ sia anche un omeomorfismo; la mappa $e^{-1}:V\iso \R^n$ definisce dunque una carta \bemph{globale} $(V,e^{-1})$ e quindi definisce da sola un atlante massimale.\\
    Banalmente, una mappa lineare $L:V\to V'$ definisce un morfismo di varietà lisce in quanto le mappe lineari $\R^n\to\R^{n'}$ sono lisce, ma non tutti i morfismi $f$ di varietà lisce sono dati da mappe lineari, dato che si richiede semplicemente che $e'\circ f\circ e^{-1}$ sia una mappa liscia, non necessariamente lineare.
\end{example}

L'esempio si generalizza immediatamente agli spazi affini di dimensione finita con lo stesso \textit{caveat}, ma il concetto di varietà è molto più generale: esempi di varietà lisce sono anche le sfere $\mathbb{S}^n$, le varietà algebriche non singolari, gli spazi proiettivi, le immagini di aperti di $\R^n$ attraverso mappe lisce e i loro grafici e moltissimi altri.

\subsubsection{Teoria delle categorie}

Nel linguaggio della teoria delle categorie, per il quale rimandiamo il lettore a \cite{Riehl2017}, possiamo esprimere elegantemente quanto detto sopra: per ogni $k \in \N_{>0}$ la seguente catena di funtori dimenticanti\footnote{Ovvero fedeli, intuitivamente un analogo dell'iniettività per i funtori.} che esibisce ciascuna di queste come categoria concreta\footnote{Moralmente, una categoria i cui oggetti sono insiemi dotati di una qualche struttura e i cui morfismi sono funzioni insiemistiche che rispettano questa struttura}:
\[ \catname{finVec}(\R) \to \catname{finAff} \to \C^\infty\catname{Man} \to \topcat \to \setcat \]
Attraverso la quale uno spazio vettoriale (dotato di una struttura di spazio affine, la topologia ereditata da $\R^n$ e la sua struttura liscia) $\left( V,0,+,\tau_e,\left\{(V, e^{-1})\right\}\right)$ perde progressivamente "pezzi" di struttura, così come le mappe che lo coinvolgono come dominio o codominio.

\newpage
\section{Geometria della fisica classica}

\newcommand\naif{\textit{naïf}\ }
Nella sezione precedente abbiamo presentato un modello molto \naif dello spaziotempo tramite spazi affini $A^1$ e $A^3$ di dimensioni $1$ e $3$, ma risulta essere troppo semplicistica per le nostre necessità, presentando due problemi particolari:\begin{itemize}
    \item Come si possono determinare le durate temporali oppure distanze e angoli nello spazio?
    \item In che modo spazio e tempo interagiscono tra loro per dar luogo a un concetto di spaziotempo "quadridimensionale"?
\end{itemize}
In questa sezione tratteremo tali questioni e svilupperemo soluzioni appropriate. 

\subsection{Spazi euclidei}

Cominciando dal primo punto, abbiamo bisogno di un "metro" sui nostri spazi per definire queste grandezze:

\begin{definition}{Prodotto scalare e norma indotta}{prodotto scalare}
    Un \bemph{prodotto scalare} su uno spazio vettoriale $V$ è una mappa $s:V\times V\to \R$ dotata delle seguenti proprietà:\begin{itemize}
        \item La bilinearità, ovvero per ogni $\lambda,\mu \in \R$ e ogni $u,v,w \in V$ si hanno 
            \[ s(\lambda u + \mu v, w ) = \lambda s(u, w) + \mu s(v, w)\qquad\text{e}\qquad s(u, \lambda v + \mu w) = \lambda s (u,v) + \mu s(u, w)\]
        \item La simmetria, ovvero $s(u,v)=s(v,u)$ per ogni $u,v \in V$;
        \item $s$ è definita positiva, ovvero $s(v,v)\ge 0$ ed in particolare si ha l'uguaglianza se e solo se $v=0$.
    \end{itemize}
    Un prodotto scalare induce una \bemph{norma} $\| \cdot \|$ su $V$ data da
    \[ \|v\| := \sqrt{s(v,v)} \]
    Che permette di associare una lunghezza $\|v\|\ge 0$ ad un vettore $v$.\\
    L'\bemph{angolo} $\vartheta \in [0,\pi]$ tra due vettori $v,w$ è definito dall'identità
    \[ s(v,w) = \|v\|\cdot \|w\|\cdot \cos(\vartheta) \]
\end{definition}

Dotare uno spazio di un prodotto scalare permette di definire una nuova categoria di spazi:

\begin{definition}{Spazio euclideo}{spazio euclideo}
    Si dice \bemph{spazio euclideo} uno spazio affine $E$ il cui spazio vettoriale sottostante $V$ è dotato di un prodotto scalare $s$.\\
    Un morfismo di spazi euclidei $\psi:E\to E'$, detto \bemph{isometria affine}, è una mappa affine la cui parte lineare $\di\psi$ preserva i prodotti scalari, ovvero 
    \[ s'(\di\psi(v),\di\psi(w)) \equiv s(v,w) \].\\
    Indichiamo con $\catname{finEuc}$ la categoria i cui oggetti sono gli spazi euclidei di dimensione finita e i cui morfismi sono le isometrie affini.
\end{definition}

Notiamo che su ogni spazio euclideo $(E,s)$ si può definire una \bemph{metrica}, o \bemph{distanza}, naturale $d: E\times E\to \R$ data da 
\[ d(P,Q) := \|P-Q\| = \sqrt{ s(P-Q, P-Q) }\]
Immediatamente si verificano gli assiomi di spazio metrico su $(E,d)$, ovvero simmetria, positività per punti distinti e disuguaglianza triangolare; inoltre, la distanza tra due punti è invariante per l'azione di $V$, cioè 
\[d(P+v,Q+v)\equiv d(P,Q)\]
Notiamo anche che una mappa affine $\psi : E\to E'$ è un'isometria se e solo se preserva le distanze, ovvero 
\[d(P,Q) \equiv d'(\psi(P),\psi(Q))\]

\renewcommand\O{\operatorname{O}}
\newcommand\SO{\operatorname{SO}}
\begin{example}{Lo spazio euclideo $\R^n$}{}
    Per $n \in \N$, lo spazio vettoriale $\R^n$ è dotato di un prodotto scalare naturale $\cdot$ definito da
    \[ \bar{v}\cdot\bar{w} = \begin{pmatrix} v^1\\\vdots\\v^n \end{pmatrix}\cdot \begin{pmatrix} w^1\\\vdots\\w^n \end{pmatrix} = \sum_{i=1}^n v^i w^i \]
    Come visto sopra, questo spazio ha anche una struttura di spazio affine, sul quale questo prodotto scalare definisce la tipica distanza 
    \[ d(\bar{x},\bar{y}) = d\left(\begin{pmatrix} x^1\\\vdots\\x^n \end{pmatrix}, \begin{pmatrix} y^1\\\vdots\\y^n \end{pmatrix}\right) = \sqrt{ \sum_{i=1}^n (x^i-y^i)^2 } \]
    Ogni isometria affine $\psi : \R^m\to\R^n$ è rappresentata tramite l'espressione $\psi(\bar{x}) = \bar{b}+L\bar{x}$ da un vettore $\bar{b\in\R^n}$ e una matrice $L$ con $n$ righe e $m$ colonne tale che $L^T L = \bar{1}_{m\times m}$.\\
    Osserviamo che tali matrici esistono solo per $n\ge m$, fatto che corrisponde alla nostra intuizione geometrica: non è possibile che $\psi$ preservi le distanze se ci sono dimensioni che vengono "schiacciate" quando questa trasforma lo spazio.\\
    Notiamo anche che quando $n=m$, abbiamo che la\footnote{mappa rappresentata dalla} matrice $L$ deve appartenere al gruppo ortogonale $\\O(n)$ e che ogni isometria $\psi:\R^n\to\R^n$ è invertibile con inversa $\psi^{-1}$ data da 
    \[ \psi^{-1}(\bar{x}) = L^T(\bar{x}-\bar{b}) = -L^T\bar{b} + L^T\bar{x} \]
\end{example}

Come per il caso vettoriale e per il caso affine, non sarà mica che ogni spazio euclideo di dimensione finita è segretamente $\R^n$ sotto mentite spoglie? E certo che è così, lo vediamo con la prossima proposizione:

\begin{proposition}{Coordinate cartesiane ortonormali}{coordinate cartesiane ortonormali}
    Sia $E$ uno spazio euclideo di dimensione $n \in \N$. Allora esiste un\footnote{non unico} isomorfismo di spazi euclidei $\psi:\R^n\iso E$, detto \bemph{sistema di coordinate cartesiane ortonormali} per $E$.
    \proof 
    Scegliamo un punto $P_0 \in E$ e una base ortonormale $\{e_i \in V : s(e_i,e_j) = \delta_{i,j}\}_{1\le i\le n}$ per lo spazio vettoriale sottostante $V$, dunque definiamo la mappa $\psi$ come
    \[\psi( (x^i)_i ) = P_0 + \sum_{i=1}^n x^ie_i \]
    Questa è ovviamente una mappa affine, notiamo che preserva i prodotti scalari:
    \[ s(\di \psi(\bar{v}),\di\psi(\bar{w}) ) = \sum_{i,j = 1}^n v^i w^j \underbrace{s(e_i,e_j)}_{=\delta_{i,j}} = \sum_{i=1}^n v^i w^i = \bar{v}\cdot\bar{w}\]
    Dunque è un'isometria. Vediamo che la sua inversa $\psi^{-1}:E\iso\R^n$ assume la forma 
    \[ \psi^{-1}(P) = \begin{pmatrix}s(e_1,P-P_0) \\ \vdots \\ s(e_n, P-P_0) \end{pmatrix} \]
    \qed
\end{proposition}

\newcommand\Aut{\operatorname{Aut}}
Grazie a questi isomorfismi è facile determinare il gruppo di automorfismi $\Aut(E)$ di uno spazio euclideo $n$-dimensionale $E$: un isomorfismo $\psi:R^n\iso E$ determina un isomorfismo $\eta_\psi : \Aut(E)\iso\Aut(\R^n)$ definito dal diagramma
\[\begin{tikzcd}
	& E && E & \\
	\\
	{\R^n} &&&& {\R^n}
	\arrow["f"{description}, from=1-2, to=1-4]
	\arrow["\psi"{description}, from=3-1, to=1-2]
	\arrow["{\eta_\psi(f) :=  \psi^{-1} \circ f \circ \psi}"{description}, from=3-1, to=3-5]
	\arrow["\psi"{description}, from=3-5, to=1-4]
\end{tikzcd}\]
Come visto sopra, abbiamo anche un isomorfismo $\R^n\times\\O(n)\iso\Aut(\R^n)$ che manda la coppia $(\bar{b},L)$ nell'automorfismo $f_{(\bar{b},L)}(\bar{x}) = \bar{b}+L\bar{x}$: i seguenti calcoli esibiscono la struttura di gruppo di $\R^n\times\\O(n)$:\begin{itemize}
    \item La composizione è data da
        \[ f_{(\bar{b},L)} \circ f_{(\bar{b}',L')(\bar{x})} = \bar{b} + L f_{(\bar{b}',L')}(\bar{x}) = \bar{b}+L\bar{b}' + LL'\bar{x} = f_{(\bar{b} + L\bar{b}', LL' )}(\bar{x}) \]
        Dunque l'operazione del gruppo è $(\bar{b},L)\cdot (\bar{b}',L') := (\bar{b} + L\bar{b}', LL' )$ e notiamo che è diversa dall'operazione \naif sul prodotto diretto del gruppo additivo $\R^n$ e del gruppo di composizione $\\O(n)$. 
    \item L'elemento neutro è ovviamente $\id = f_{(\bar{0},\bar{1}_{n\times n})}$.
    \item L'inverso di $(\bar{b},L)$ è dato da $(-L^T\bar{b},L^T)$.
\end{itemize}
Queste trasformazioni hanno il significato geometrico di traslazioni, rotazioni e riflessioni dello spazio euclideo $\R^n$.\\
Applicando subito quanto visto al tempo $E^1$ e allo spazio $E^3$ della fisica classica otteniamo in qualunque scelta di coordinate una descrizione esplicita dei loro automorfismi: 
\[ \R^1\iso \R^1, t\mapsto at+c \]
\[ \R^3 \iso \R^3, \bar{x}\mapsto L\bar{x} + \bar{b} \]
Per $a \in \{-1,1\}$, $c \in \R$, $L\in\\O(3)$ e $\bar{b}\in\R^3$; per escludere trasformazioni che possano riflettere il tempo (con $a=-1$) o lo spazio (con $\det L = -1$), dotiamo questi nostri spazi euclidei di un'\bemph{orientazione} ereditata da quella sul loro spazio vettoriale sottostante e richiediamo che gli automorfismi la preservino:

\begin{definition}{Orientazione}{}
    Un'\bemph{orientazione} su uno spazio vettoriale $V$ è una classe di equivalenza di basi $[\{e_i\}_i]_\sim$ dove due basi sono equivalenti se e solo se la mappa di cambiamento di base $f(e_i') = e_i$ ha determinante positivo.
\end{definition}

Nel caso orientato, siamo forzati a scegliere $a=1$ e $L \in \SO(3)$, il sottogruppo del gruppo ortogonale detto \bemph{gruppo ortogonale speciale o proprio} delle matrici $3\times 3$ ortogonali con determinante uguale a $1$.

\subsection{Fibrati}

Inizialmente potremmo pensare ancora una volta a una soluzione \naif per la struttura del nostro spaziotempo unificato, ovvero il prodotto $E^1\times E^3$ del tempo e dello spazio, ma ancora una volta la realtà dei fatti decide di rovinarci la festa e ci costringe a lavorare di più: questa descrizione ingenua presumerebbe un prodotto in qualche modo "privilegiato", di cui non c'è evidenza nella natura.\\
Qualcosa che ha senso definire \textit{nel contesto della meccanica classica} è la mappa di proiezione sul tempo 
\[\pi_1: E^1\times E^3 \to E^1\] 
che manda il punto $(t,P)$ nel suo tempo $t$, che fisicamente interpretiamo come il \bemph{tempo assoluto}.\\
Seguendo \cite{Moretti2024}, formuleremo il nostro concetto di spaziotempo in modo indipendente da coordinate usando il linguaggio delle varietà, per il quale avremo bisogno di alcuni strumenti di geometria differenziale.

\begin{definition}{Sommersione}{sommersione}
    Un morfismo di varietà lisce $f:\M\to\M'$ di dimensioni $n, n' \in \N$ si dice \bemph{sommersione in $p \in \M$}\footnote{Nella terminologia di \cite{Moretti2024}, non singolare} se per ogni\footnote{E dunque per una qualunque} scelta di carte locali $(U,\phi)$ e $(U',\phi')$ con $f(U)\subset U'$ e $p\in U$, la matrice Jacobiana
    \[ \left[ \frac{\partial (\phi\circ f\circ \phi^{-1})^j  }{\partial x^i} \bigg|_{\phi(p)} \right]_{1\le i\le n, 1\le j\le n'} : \R^n \to \R^{n'}\]
    della mappa liscia $\phi\circ f\circ\phi^{-1}$ ha rango massimo.\\
    Un morfismo si dice \bemph{sommersione} se lo è per ogni punto $p\in\M$.
\end{definition}

Ed ora, finalmente

\begin{definition}{Spaziotempo della fisica classica}{spaziotempo}
    Uno \bemph{spaziotempo della fisica classica} è una varietà liscia $\V^4$ di dimensione $4$ dotata delle seguenti strutture:\begin{itemize}
        \item Una sommersione suriettiva $T:\V^4\to E^1$ detta \bemph{tempo assoluto};
        \item Per ogni $t \in E^1$, una struttura di spazio euclideo di dimensione $3$ sulla fibra $\Sigma_t := T^{-1}(\{t\})$ detta \bemph{spazio assoluto al tempo} $t$.
    \end{itemize} 
    Queste strutture formano un \bemph{fibrato} $T:\V^4\to E^1$, ovvero sono soggette alla seguente condizione di compatibilità: per ogni $t \in E^1$ deve esistere un intorno aperto $I$ di $t$ e un diffeomorfismo\footnote{isomorfismo di varietà lisce} $F: T^{-1}(I)\iso I\times E^3$ tale che:\begin{itemize}
        \item Il seguente diagramma commuti 
            \[\begin{tikzcd}
        	    {T^{-1}(I)} &&&& {I\times E^3} \\
	            && I
	            \arrow["F"{description}, from=1-1, to=1-5]
	            \arrow["T"{description}, from=1-1, to=2-3]
	            \arrow["{\pi_I}"{description}, from=1-5, to=2-3]
            \end{tikzcd}\]
        \item Per ogni $t'\in I$, la restrizione $F|_{\Sigma_{t'}} : \Sigma_{t'}\iso \{t'\}\times\Sigma_{t'}\cong \Sigma_{t'}$ sia un'isometria affine.
    \end{itemize}
    Uno spaziotempo viene detto \bemph{orientato} se tutte le fibre sono dotate di orientazioni tali che ogni $F|_{\Sigma_{t'}}$ le preservi.\\
    I punti di $\V^4$ si dicono \bemph{eventi}, e rappresentano fisicamente il quando e il dove avviene qualcosa.
\end{definition}

Si noti però che abbiamo scritto \bemph{uno}\footnote{Il lettore esasperato potrebbe essere comprensibilmente tentato dal lanciare un grido di frustrazione e decidere di dedicarsi a passatempi più edificanti, come il curling o la transumanza sull'Appennino abruzzese, ma lo assicuriamo che ormai siamo in dirittura d'arrivo} spaziotempo, in linea di principio potrebbero essercene uno come dieci, cento, mille, numerabili, qualche cardinale inaccessibile, una classe propria... Ma vengono in nostro aiuto degli strumenti più avanzati.

\begin{definition}{Isomorfismo di spaziotempi}{}
    Un \bemph{isomorfismo} tra due spaziotempi $T:\V^4\to E^1$ e $T':\V'^4 \to E'^1$ è una coppia $(F,f)$ composta da un diffeomorfismo $F:\V^4\iso\V'^4$ e un'isometria affine $f: E^1\iso E'^1$ tale che il seguente diagramma commuti 
    \[\begin{tikzcd}
    	{\V^4} && {\V'^4} \\
	    \\
	    {E^1} && {E'^1}
	    \arrow["F"{description}, from=1-1, to=1-3]
	    \arrow["T"{description}, from=1-1, to=3-1]
	    \arrow["{T'}"{description}, from=1-3, to=3-3]
	    \arrow["f"{description}, from=3-1, to=3-3]
    \end{tikzcd}\]
    E che per ogni $t \in E^1$ la restrizione alle fibre $F|_{\Sigma_t} :\Sigma_t\iso \Sigma'_{f(t)}$ sia un'isometria affine.
\end{definition}

Enunciamo ora un teorema di cui non possiamo dare la dimostrazione perchè richiede strumenti più avanzati di geometria differenziale che non possediamo e che esulano dagli scopi di questo corso.

\begin{theorem}{Essenziale unicità dello spaziotempo della fisica classica}{spaziotempo standard}
    Ogni spaziotempo della fisica classica $T:\V^4\to E^1$ è isomorfo in modo non canonico allo spaziotempo standard $\pi_1: E^1 \times E^3\to E^1$.\\
    Una scelta di un isomorfismo 
    \[ (\mc R, \id_{E^1}) : (T:\V^4\to E^1) \iso (\pi_1: E^1 \times E^3\to E^1)  \]
    È detta \bemph{(sistema di) riferimento}.
\end{theorem}

Un cambiamento di riferimento dunque è dato da un'automorfismo dello spaziotempo standard
\[\begin{tikzcd}
	& {T:\V^4\to E^1} & \\
	\\
	{\pi_1:E^1\times E^3 \to E^1} && {\pi_1: E^1 \times E^3\to E^1}
	\arrow["{\mc R}"{description}, from=1-2, to=3-1]
	\arrow["{\mc R'}"{description}, from=1-2, to=3-3]
	\arrow["{\mc R' \circ \mc R^{-1}}"{description}, from=3-1, to=3-3]
\end{tikzcd}\]
Che di seguito descriveremo esplicitamente.\\
Come abbiamo visto sopra, uno spazio euclideo $E^n$ ammette una scelta (non canonica) di coordinate data da $\psi_n:\R^n\to E^n$, dunque possiamo ottenere una descrizione ancora più esplicita dello spaziotempo attraverso delle \bemph{coordinate standard adattate al riferimento}
\[\begin{tikzcd}
	{\V^4} && {E^1\times E^3} && {\R^1\times \R^3} \\
	\\
	{E^1} && {E^1} && {\R^1}
	\arrow["{\mc R}"{description}, from=1-1, to=1-3]
	\arrow["T"{description}, from=1-1, to=3-1]
	\arrow["{(\psi_1^{-1},\psi_3^{-1})}"{description}, from=1-3, to=1-5]
	\arrow["{\pi_1}"{description}, from=1-3, to=3-3]
	\arrow["{\pi_1}"{description}, from=1-5, to=3-5]
	\arrow[equals, from=3-1, to=3-3]
	\arrow["{\psi_1^{-1}}"{description}, from=3-3, to=3-5]
\end{tikzcd}\]
La potenza delle coordinate è che ci permettono di effettuare calcoli: vediamo un esempio:  abbiamo un isomorfismo di gruppi
\[ \Aut(T:\V^4\to E^1) \iso \Aut(\pi_1:\R^1\times\R^3\to\R^1 ) \]
E notiamo che il secondo è notevolmente più semplice da determinare del primo

\begin{proposition}{Gruppo degli automorfismi di $\pi_1:\R^1\times\R^3\to\R^1$}{automorfismi dello spaziotempo standard}
    Ogni automorfismo 
    \[ (F,f): (\pi_1:\R^1\times\R^3\to\R^1)\iso (\pi_1:\R^1\times\R^3\to\R^1) \]
    È della forma 
    \[ F: (t,\bar{x}) \mapsto (c+at, \bar{b}(t) + R(t)\bar{x} ) \quad f:t\mapsto c+at \]
    Con $c$ in $\R$, $a =\pm 1$, $\bar{b}$ in $\C^{\infty}(\R^1,\R^3)$ e $R$ in $\C^{\infty}(\R^1,\O(3))$.\\
    L'insieme sottostante al gruppo degli automorfismi allora è dato dal prodotto cartesiano
    \[ \R^1\times \{\pm 1\} \times \C^{\infty}(\R^1,\R^3) \times \C^{\infty}(\R^1,\O(3)) \]
    E la struttura di gruppo è data da
    \[ (c,a,\bar{b},R)(c', a', \bar{b}', R') = (c+ac', aa', \bar{b}(c'+a'(\cdot))+R(c'+a'(\cdot))\bar{b}'(\cdot) , R(c'+a'(\cdot))R'(\cdot) ) \]
    con evidente neutro $(0,1, \bar{0},\id_{\R^3})$ e inversi dati da
    \[ (c,a,\bar{b}, R)^{-1} = (-ac, a, -R^T(-ac+a(\cdot))\bar{b}(-ac+a\cdot), R^T(-ac+a(\cdot)) ) \]
    \proof 
    Abbiamo già visto che le isometrie affini di $\R^1$ sono della forma $t\mapsto c+at$; per la proprietà universale del prodotto, il diffeomorfismo $F$ di $\R^1\times\R^3$ si può scrivere in componenti $(F_1, F_3)$, dunque notiamo la commutatività del seguente diagramma 
    \[\begin{tikzcd}
    	{\R^1\times \R^3} &&&& {\R^1\times\R^3} \\
	    & {(t,\bar{x})} && {(F_1(t,\bar{x}),F_3(t,\bar{x}))} \\
	    \\
	    & t && {f(t) = F_{1}(t,\bar{x})} \\
	    {\R^1} &&&& {\R^1}
	    \arrow["F"{description}, from=1-1, to=1-5]
	    \arrow["{\pi_1}"{description}, from=1-1, to=5-1]
	    \arrow["{\pi_1}"{description}, from=1-5, to=5-5]
	    \arrow[maps to, from=2-2, to=2-4]
	    \arrow[maps to, from=2-2, to=4-2]
	    \arrow[maps to, from=2-4, to=4-4]
	    \arrow[maps to, from=4-2, to=4-4]
	    \arrow["f"{description}, from=5-1, to=5-5]
    \end{tikzcd}\]
    Che ci da $F(t,\bar{x}) = c+at$; dato che richiediamo che $\bar{x}\mapsto F_3(t,\bar{x})$ sia un'isometria affine, deve essere della forma $(t,\bar{x})\mapsto  \bar{b}(t)+ R(t)\bar{x}$ con $\bar{b}(t)$ in $\R^3$ e $R(t)$ in $\O(3)$, ma dato che deve essere differenziabile queste quantità devono dipendere in modo differenziabile da $t$, dunque ecco dimostrata la biezione; la verifica delle leggi di gruppo è banale e lasciata come esercizio al lettore.
    \qed
\end{proposition}

Nel caso in cui lo spaziotempo dovesse essere dotato di un'orientazione, la proposizione di cui sopra varrebbe ugualmente limitandosi ai sottogruppi $\{1\}$ e $\SO(3)$ rispettivamente di $\{\pm 1\}$ e $\O(3)$.\\
In particolare, se lo spaziotempo dovesse presentare una struttura di spazio euclideo, un automorfismo preserverebbe questa struttura se e solo se fosse della forma
\[ (t,\bar{x})\mapsto (c+at, \bar{b} + \bar{v}t + R\bar{x}) \]
con $c,a,\bar{b},\bar{v}$ e $R$ costanti\footnote{adeguatamente ristrette nel caso orientato}; questo caso descrive il gruppo delle \bemph{trasformazioni Galileiane} e più avanti vedremo come le leggi di Newton inducano una struttra affine.

\chapter{Meccanica Newtoniana}

\section{Cinematica}

L'oggetto fondamentale di interesse della meccanica classica è il \bemph{punto materiale}, o \bemph{particella}, che si muove nello spazio nel corso del tempo. 

\begin{definition}{Linea di universo}{linea di universo}
    Una \bemph{linea di universo} nello spaziotempo $T:\V^4\to E^1$ è una curva liscia $\gamma : I \to \V^4$ da un intervallo aperto $I\subset E^1$ tale che il seguente diagramma commuti 
    \[\begin{tikzcd}
	    && {\V^4} \\
	    \\
	    I && {E^1}
	    \arrow["T"{description}, from=1-3, to=3-3]
	    \arrow["\gamma"{description}, from=3-1, to=1-3]
	    \arrow[hook, from=3-1, to=3-3]
    \end{tikzcd}\]
    Un punto materiale verrà identificato con la sua linea di universo.
\end{definition}

Notiamo che le linee di universo si comportano bene sotto l'azione degli isomorfismi di spaziotempi, infatti per $\gamma$ e $T$ come nella definizione e un isomorfismo $(F,f) : T\to T'$ vediamo facilmente che la mappa $\gamma ' := F\circ \gamma \circ f^{-1}$ è liscia soddisfa il requisito $T'\circ \gamma' = \id_I$, quindi è una linea di universo in $T'$.\\
Specializzandoci con un riferimento al caso di $\pi_1:E^1\times E^3\to E^1$, notiamo che una linea di universo può essere scritta in componenti come $(\gamma_1,\gamma_3)$ e che il requisito $\pi_1\circ \gamma = \id_I$ implica che $\gamma_1$ sia l'identità; di conseguenza, le linee di universo nello spaziotempo standard (e dunque di qualsiasi spaziotempo) sono in corrispondenza biunivoca con le curve lisce in $E^3$.\\
Specializzandoci ulteriormente con una scelta di coordinate, vediamo che una linea di universo $\gamma$ è determinata da una curva liscia $\bar{\gamma}$ in $\R^3$; adottando questo punto di vista, l'azione degli automorfismi calcolati in \cref{prop:automorfismi dello spaziotempo standard} sulle linee di universo può essere calcolata esplicitamente: dato un automorfismo $(F,f)$ con $f(t)= c+at$ e $F(t,\bar x)= (c+at, \bar b(t) + R(t)\bar x )$ , vediamo che $F\circ\gamma\circ f^{-1}$ è dato dalla mappa $t' \mapsto (t', \bar b(-ac +at') + R(-ac+at')\bar \gamma (-ac +at') )$ che ci permette di trasformare $\bar \gamma : I \to \R^3$ in $\bar \gamma' : aI+c \to \R^3$.

\newcommand\tv{\text{(tv)}}
\newcommand\cor{\text{(cor)}}
\subsection{Velocità e accelerazione}

Possiamo finalmente definire le due quantità centrali nella cinematica, ovvero la velocità e l'accelerazione. A priori tuttavia l'intuizione della velocità come derivata temporale non è adatta al nostro ambiente, poichè per tempi diversi $t$ e $t+h$ si ha che $\gamma$ appartiene a spazi euclidei $\Sigma_t$ e $\Sigma_{t+h}$ che sono a priori distinti, dunque non è possibile calcolare il rapporto incrementale in questo modo; definiamo quindi un concetto di derivata rispetto a un riferimento $R$ e poi studieremo come si comporta sotto l'azione dei cambiamenti di coordinate.

\begin{definition}{Derivata rispetto a un riferimento}{derivata rispetto a un riferimento}
    Siano $T:\V^4\to E^1$ uno spaziotempo, $\gamma : I\to\V^4$ una linea di universo e $R:T\to\pi_1$ un riferimento. Definiamo la mappa $\gamma_R : I \to E^3$ data da $t\mapsto R|_{\Sigma_t}(\gamma(t))$. Si dice \bemph{derivata temporale rispetto al riferimento} al tempo $t$ l'elemento di $V^3$ spazio vettoriale sottostante a $E^3$ dato dal limite 
    \[ \frac{\di}{\di t}\gamma_R(t) := \lim_{h \to 0}\frac{\gamma_R(t+h)-\gamma_R(t)}{h} \]
    Utilizzando la parte lineare $\di R|_{\Sigma_t} : V_t\iso V^3$ si ottiene un elemento di $V_t$ e dunque definiamo
    \[ \frac{\di}{\di t}\bigg|_R \gamma(t) := \left( \di R|_{\Sigma_t}^{-1} \circ \frac{\di}{\di t} \circ R|_{\Sigma_t} \right)(\gamma(t)) \]
\end{definition}

Possiamo dunque definire in modo adatto velocità e accelerazione:

\begin{definition}{Velocità e accelerazione rispetto a un riferimento}{velocità e accelerazione}
    Siano $T:\V^4\to E^1$ uno spaziotempo, $\gamma : I\to\V^4$ una linea di universo e $R:T\to\pi_1$ un riferimento.\\
    Definiamo la \bemph{velocità} di $\gamma$ rispetto a $R$ come 
    \[ v|_R(t) := \frac{\di}{\di t}\bigg|_R \gamma(t) = \left( \di R|_{\Sigma_t}^{-1} \circ \frac{\di}{\di t} \circ R|_{\Sigma_t} \right)(\gamma(t)) \]
    E definiamo la sua \bemph{accelerazione} come 
    \[ a|_R(t) := \frac{\di^2}{\di t^2}\bigg|_R \gamma(t) = \left( \di R|_{\Sigma_t}^{-1} \circ \frac{\di^2}{\di t^2} \circ R|_{\Sigma_t} \right)(\gamma(t)) \]
\end{definition}

Sarà utile esprimere queste formule in una qualsiasi scelta di coordinate standard adattate a $R$, ovvero una scelta di un punto $O \in E^3$ e una base ortonormale $\{e_1,e_2,e_3\}\subset V^3$: definendo
\[ O(t):= R|_{\Sigma_t}^{-1}(O) \in \Sigma_t \]
\[ e_i(t) := \di R|_{\Sigma t}^{-1}(e_i) \in V_t \]
Possiamo scrivere 
\[ \gamma(t) = O(t) + \sum_i \bar\gamma^i(t) e_i(t) \in \Sigma_t \]
E dunque ottenere
\begin{align*} 
    v|_R(t) & = \di R|_{\Sigma_t}^{-1}\ \frac{\di}{\di t}\left( R|_{\Sigma_t}(O(t)) + \sum_i \bar\gamma^i(t) \di R|_{\Sigma_t}(e_i(t)) \right) \\ & = \di R|_{\Sigma_t}^{-1}\ \frac{\di}{\di t}\left( O + \sum_i\bar\gamma^i(t) e_i \right) \\ & = \sum_i \frac{\di}{\di t} \bar\gamma^i(t) \di R|_{\Sigma_t}^{-1}(e_i) = \sum_i \frac{\di}{\di t} \bar\gamma^i(t)e_i(t)
\end{align*}
E in modo del tutto analogo per l'accelerazione 
\[ a|_R(t) = \sum_i \frac{\di^2}{\di t^2} \bar\gamma^i(t)e_i(t) \]
Avendo definito velocità e accelerazione rispetto a un riferimento $R$, è importante capire come queste cambino sotto il passaggio a un altro sistema di riferimento $R'$.\\
Scegliendo due sistemi di coordinate adattati $(O,\{e_i\}_i)$ e $(O', \{e'_i\}_i)$, possiamo scrivere una linea di universo $\gamma(t) \in \Sigma_t$ in due modi diversi:
\[ \gamma(t) = O(t) + \sum_i \bar\gamma^i(t) e_i(t) = O'(t) + \sum_i \bar\gamma'^i(t) e'_i(t) \]
Per confrontare $v|_R$ e $v|_{R'}$, usiamo che esiste un $A \in \C^\infty(I, \O(3))$ tale che 
\[ e_i(t) = \sum_j A_i^j(t) e'_j(t) \]
e calcoliamo 
\begin{align*} \gamma(t) & = O'(t) + \left( O(t)-O'(t) +\sum_{i,j} \bar\gamma^i(t) A_i^j(t) e'_j(t)  \right)\\ & = O'(t) + \sum_j\underbrace{\left( b^j(t) + \sum_i A_i^j(t)\bar\gamma^i(t) \right)}_{= \bar\gamma'^j(t)} e'_j(t) \end{align*}
dove $\bar b \in \C^{\infty(I, \R^3)}$ è data da $O(t)-O'(t) = \sum_i b^i(t)e'_i(t)$: otteniamo la relazione
\begin{align*}
    v|_{R'}(t) & = \sum_i \frac{\di}{\di t}\bar\gamma'^i(t) e'_i(t) \\ & = \sum_i\frac{\di}{\di t} \left( b^i(t) + \sum_j A^i_j(t)\bar\gamma^j(t) \right)e'_i(t) \\ & = \sum_i\left[ \frac{\di}{\di t} b^i(t) + \sum_j\left( \frac{\di A^i_j(t) }{\di t}\bar\gamma^j(t) + A^i_j(t) \frac{\di \bar\gamma^j(t)}{\di t} \right) \right] e'_i(t) \\ & = \underbrace{\sum_j \frac{\di \bar\gamma^j(t)}{\di t} e_j(t)}_{=v|_R(t)} + \underbrace{\sum_i \left[ \frac{\di b^i(t)}{\di t} + \sum_j\left( \frac{\di A^i_j(t)}{\di t} \bar\gamma^j(t) \right) \right]e'_i(t)}_{=: v^\tv (t)}
\end{align*} 
dove $v^\tv(t)$ è detta \bemph{velocità di trascinamento} di $R$ rispetto a $R'$. Con un calcolo analogo troviamo che l'accelerazione ha una forma simile:
\[ a|_{R'}(t) = \underbrace{\sum_j \frac{\di^2 \bar\gamma^j(t)}{\di t^2} e_j(t)}_{=a|_R(t)} + \underbrace{\sum_i \left[ \frac{\di^2 b^i(t)}{\di t^2} + \sum_j\left( \frac{\di^2 A^i_j(t)}{\di t^2} \bar\gamma^j(t) \right) \right]e'_i(t)}_{=: a^\tv (t)} + \underbrace{2\sum_{i,j} \frac{\di A^i_j(t)}{\di t} \frac{\di \bar\gamma^j(t)}{\di t} e'_i(t)}_{=: a^\cor(t)} \]
dove $a^\tv(t)$ è detta \bemph{accelerazione di trascinamento} di $R$ rispetto a $R'$ e $a^\cor(t)$ è detta \bemph{accelerazione di Coriolis} e caratteristicamente dipende dalla velocità $v|_R(t)$.

\newpage
\section{Dinamica Newtoniana}

In questa sezione, discutiamo i principi della dinamica di Newton, che governano la dinamica di (potenzialmente molteplici) punti materiali e hanno riscosso grandissimo successo nel descrivere situazioni reali (con le dovute approssimazioni, ovviamente).

\subsection{Primo principio, o principio d'inerzia}

Il primo principio stabilisce una classe privilegiata di sistemi di riferimento per lo spaziotempo $T:\V^4\to E^1$, caratterizzati dal comportamento di sistemi di particelle "isolati", ovvero separate tra loro e dagli altri corpi dell'universo da grandissime distanze.

\begin{axiom}{Primo principio della dinamica}{dinamica 1}
    Esiste una classe di sistemi di riferimento detti \bemph{inerziali} in ciascuno dei quali ogni punto di un insieme arbitrario di punti materiali isolati si muove di moto rettilineo uniforme, ovvero a velocità costante dipendente soltanto dal punto materiale considerato.
\end{axiom}

\textit{Espressiamo} (sic) questo principio in un linguaggio vagamente più formale.\\
Un insieme di $N$ punti materiali $\{ Q_1,\ ...\ ,\ Q_N \}$ corrisponde a $N$ linee di universo $\gamma_{Q_k}:I\to\V^4$ nello spaziotempo. Dato in qualsiasi sistema di riferimento $(R,\id_{E^1})$ per lo spaziotempo con coordinate standard adattate, abbiamo introdotto il concetto di velocità rispetto ad $R$
\[ v_{Q_k}|_R (t) = \sum_i\frac{\di \bar\gamma_{Q_k}^i(t)}{\di t}e_i(t) \]
Il primo principio allora richiede che l'accelerazione di ciascun punto sia nulla, ovvero
\[ a_{Q_k}|_R (t) = \sum_i\frac{\di^2 \bar\gamma_{Q_k}^i(t)}{\di t^2}e_i(t) = 0 \]
È anche importante sottolineare che non c'è un solo sistema di riferimento inerziale, ma ce ne sono a priori quanti vogliamo: ricordando che l'automorfismo $(F,f)$ nel diagramma 
\[\begin{tikzcd}
	&& {T:\V^4 \to E^1} && \\
	{\pi_1:E^1\times E^3 \to E^1} &&&& {\pi_1:E^1\times E^3 \to E^1} \\
	\\
	{\pi_1:\R^1\times\R^3\to\R^1} &&&& {\pi_1:\R^1\times\R^3 \to \R^1}
	\arrow["{(R,\id_{E^1})}"{description}, from=1-3, to=2-1]
	\arrow["{(R',\id_{E^1})}"{description}, from=1-3, to=2-5]
	\arrow[from=2-1, to=4-1]
	\arrow[from=2-5, to=4-5]
	\arrow["{(F,f)}"{description}, from=4-1, to=4-5]
\end{tikzcd}\]
Ha forma generale $f(t) = c+at$ e $F(t,\bar x) = (c+at, \bar b(t) + A(t) \bar x)$ con $c \in \R$, $a \in \{\pm1\}$, $\bar b \in \C^{\infty}(\R^1, \R^3)$ e $A \in \C^{\infty}(\R^1, \O(3))$, il seguente risultato è di importanza fondamentale:

\begin{theorem}{Principio di inerzia e trasformazioni di Galileo}{inerzia e Galileo}
    Assumiamo che per ogni riferimento $R$ e ogni evento $p\in\V^4$ esista un punto materiale la cui linea di universo attraversi $p$ con velocità (rispetto a $R$) fissabile a piacimento e che il punto sia isolato in qualche intervallo aperto $I \subset E^1$ con $T(p)\in I$\\
    Se $R$ è inerziale, dato un altro riferimento $R'$ questo è inerziale se e solo se l'automorfismo $(F,f)$ che descrive il cambiamento di sistema di riferimento in coordinate standard adattate assume la forma di una trasformazione di Galileo, ovvero $f(t) = c+at$ e $F(t,\bar x) = (c+at, \bar b + \bar v t A \bar x)$ con $c \in \R$, $a \in \{\pm1\}$, $\bar b,\bar v \in \R^3$ e $A \in \O(3)$ costanti.
    \proof 
    Ricordiamoci le formule di trasformazione derivate sopra
    \[ a|_R'(t) = a|_R(t) + a^\tv(t) + a^\cor(t) \]
    Dove
    \[ a^\tv(t) = \sum_i \left[ \frac{\di^2 b^i(t)}{\di t^2} + \sum_j\frac{\di^2 A^i_j(t)}{\di t^2} \bar\gamma^j(t) \right]e'_i(t) \]
    \[ a^\cor(t) = 2\sum_{i,j} \frac{\di A^i_j(t)}{\di t} \frac{\di \bar\gamma^j(t)}{\di t} e'_i(t) \]
    E dimostriamo entrambe le implicazioni:\begin{itemize}
        \item Supponiamo che $(F,f)$ sia una trasformazione di Galileo: notiamo immediatamente che dalle formule di trasformazione, dato un insieme di punti isolati $\{Q_k\}_k$, per ciascuno di essi $a_{Q_k}|_R(t)$ è nulla per definizione, $a^\cor(t)$ è nulla per la costanza di $A$ e aggiungendo la costanza\footnote{A dire il vero, funzionerebbe anche se $\bar b(t)$ variasse solo al primo ordine} di $\bar b(t)$ otteniamo anche la nullità di $a^\tv(t)$.
        \item Supponiamo che $R$ e $R'$ siano inerziali; in particolare, per ogni singolo punto isolato $Q$ abbiamo $a_Q|_R =0$ e $a_Q|_{R'} =0$. Come conseguenza dell'ipotesi di cui sopra, ogni linea di universo 
        \[ \gamma_Q(t) = O(t) +\sum_i (q^i +c^i t )e_i(t) \]
        con $\bar q, \bar c \in \R^3$ descrive un tale punto isolato con accelerazione nulla rispetto a $R$; dato che anche l'accelerazione rispetto a $R'$ deve essere nulla, abbiamo che $a^\tv(t)+a^\cor(t) = 0$, ovvero 
        \[ 0= \sum_i\left[ \frac{\di^2 b^i(t)}{\di t^2} + \sum_j \frac{\di^2 A^i_j(t)}{\di t^2}(q^j + c^j t) + 2\sum_j \frac{\di A^i_j(t)}{\di t}c^i  \right]e'_i(t) \qquad \forall \bar q, \bar c \in \R^3\]
        Osserviamo che tre casi particolari determinano le variazioni di $\bar b$ e $A$:
        \[ \bar q = \bar c = 0 \Rarr \frac{\di^2 \bar b(t)}{\di t^2} = 0\Rarr \bar b(t) = \bar b + \bar v t\]
        \[ \bar q \in \R^3, \bar c = 0 \Rarr \frac{\di^2 A(t)}{\di t^2} = 0 \Rarr A(t) = A + K t \]
        \[\bar q = 0, \bar c \in \R^3  \Rarr \frac{\di A(t)}{\di t} = 0 \Rarr A(t) = A \]
        Quindi il cambiamento tra sistemi di riferimento è una trasformazione Galileiana.
    \end{itemize}
    \qed
\end{theorem}

Una conseguenza importante del primo principio è che questo permette di dare allo spaziotempo $T:\V^4\to E^1$ una struttura affine privilegiata 

\begin{theorem}{Struttura affine su $T:\V^4\to\R^1$}{spaziotempo affine}
    Esiste una struttura affine $+:\V^4\times V^4 \to \V^4$ sullo spaziotempo $T:\V^4\to E^1$ tale che per ogni riferimento inerziale $(R,\id_{E^1})$ la mappa $R:\V^4\to E^1\times E^4$ sia affine.\\
    Questa struttura è unica a meno di\footnote{non unici} isomorfismi affini che agiscono su $\V^4$ come l'identità ed è detta \bemph{struttura Galileiana}.
    \proof 
    Ricordando che $E^1\times E^3$ ha canonicamente una struttura affine su $V^1\times V^3$ data da $(t,O)+(c,v) = (t+c, O+v)$, per qualsiasi riferimento inerziale $R$ possiamo trasferirla a $\V^4$ tramite il diffeomoerfismo $R$ definendo $ p +_R (c,v) = R^{-1}(R(p) + (c,v))$. Considerando un qualsiasi altro riferimento inerziale $R'$, sappiamo dal teorema precedente che il cambiamento di sistema di riferimento
    \[\begin{tikzcd}
    	& {\V^4} & \\
	    \\
	    {E^1\times E^3} && {E^1\times E^3}
	    \arrow["R"{description}, from=1-2, to=3-1]
	    \arrow["{R'}"{description}, from=1-2, to=3-3]
	    \arrow["{R'\circ R^{-1}}"{description}, from=3-1, to=3-3]
    \end{tikzcd}\]
    deve assumere la forma di una trasformazione Galileiana, dunque deve essere affine rispetto alla struttura indotta da $R$ e quindi anche $R'$ deve essere una trasformazione affine.\\
    Dimostriamo l'unicità: consideriamo un'altra struttura affine $+$ su $\V^4$ che soddisfi le proprietà richieste; in particolare per quanto visto sopra, il seguente diagramma deve commutare
    \[\begin{tikzcd}
    	{\V^4_+} && {\V^4_{+_R}} \\
	    \\
	    & {E^1\times E^3}
	    \arrow["{R\circ R^{-1} = \id_{\V^4}}"{description}, from=1-1, to=1-3]
	    \arrow["R"{description}, from=1-1, to=3-2]
	    \arrow["R"{description}, from=1-3, to=3-2]
    \end{tikzcd}\]
    E dunque le due strutture affini devono essere isomorfe. Si noti che la parte lineare $V^4\iso V^1\times V^3$ di questo isomorfismo non è necessariamente l'identità però.
    \qed
\end{theorem}

È interessante notare che la mappa $T:\V^4\to E^1$ è affine, in quanto per qualsiasi riferimento inerziale $R$ vale $T = \pi_1 \circ R$.\\
La verifica che il gruppo degli automorfismi di questa struttura affine è esattamente il gruppo delle trasformazioni Galileiane è lasciata come esercizio per il lettore.

\subsection{Secondo principio e le forze}

A differenza del primo principio, questo e i prossimi permettono che i singoli punti del sistema considerato non siano necessariamente isolati tra loro. Questi vanno formulati comunque rispetto alla classe dei riferimenti inerziali e assumendo che l'intero sistema sia isolato rispetto al resto dell'universo.
\paragraph{Notazione.} Dato un riferimento inerziale $(R,\id_{E^1})$, d'ora in avanti identificheremo sempre gli spazi euclidei $\Sigma_t \iso E^3$ e gli spazi vettoriali sottostanti $V_t\iso V^3$; in particolare, considereremo sempre una linea di universo come una curva liscia $\gamma |_R :I\to E_3$ con velocità e accelerazione date semplicemente da 
\[ v|_R = \frac{\di }{\di t}\gamma|_R : I \to V^3, \qquad a|_R = \frac{\di^2 }{\di t^2}\gamma|_R : I \to V^3\]

\begin{axiom}{Principi di massa e impulso}{massa e impulso}
    Per ogni punto materiale $Q$ esiste una costante $m_Q>0$ detta \bemph{massa} di $Q$ tale che valgano i seguenti fatti:\begin{itemize}
        \item Principio di conservazione dell'impulso: in ogni riferimento inerziale e per ogni sistema isolato costituito da due punti $P,Q$, il vettore $m_P v_P|_R + m_Q v_Q|_R$ è costante nel tempo anche se non lo sono separatamente le singole velocità; il vettore $p_Q = m_Q v_Q|_R$ è detto \bemph{impulso} di $Q$ rispetto a $R$. 
        \item Principio di conservazione della massa: nei processi in cui il numero dei punti materiali cambia, la somma delle masse iniziali è uguale alla somma delle masse finali.
    \end{itemize} 
\end{axiom}

\begin{axiom}{Secondo principio della dinamica}{dinamica 2.1}
    Per ogni sistema di riferimento $R$ e ogni sistema isolato costituito da due punti isolati $Q,Q'$, esistono due mappe lisce 
    \[ F_R, F'_R : E^3\times E^3 \times V^3\times V^3 \to V^3 \times V^3 \] tali che 
    \[ \forall t \in I, \begin{dcases}
        m_Q a_Q|_R(t) = F_R(\gamma_Q|_R(t), \gamma_{Q'}|_R(t), v_Q|_R(t), v_{Q'}|_R(t))\\
        m_{Q'} a_{Q'}|_R(t) = F'_R(\gamma_{Q'}|_R(t), \gamma_{Q}|_R(t), v_{Q'}|_R(t), v_{Q}|_R(t))
    \end{dcases}\]
    $F_R$ ($F'_R$) è chiamata la \bemph{forza} che agisce su $Q$ ($Q'$) a causa di $Q'$ ($Q$), e ciascuna è detta \bemph{reazione} rispetto all'altra.\\
    Inoltre, le forze sono equivarianti rispetto a cambiamenti di sistema di riferimento inerziale, ovvero il\footnote{A dire il vero, i seguenti due diagrammi, dato che ne otteniamo un altro scambiando $F$ con $F'$} seguente diagramma commuta per ogni $t \in I$:
    \[\begin{tikzcd}
    	&& {\Sigma_t\times\Sigma_t \times V_t\times V_t} && \\
	    \\
	    {E^3\times E^3\times V^3\times V^3} && {V_t} && {E^3\times E^3\times V^3\times V^3} \\
	    \\
	    {V_3} &&&& {V^3}
	    \arrow["{R'|_{\Sigma_t}\times R'|_{\Sigma_t} \times \di R'|_{\Sigma_t} \times \di R'|_{\Sigma_t}}"{description}, from=1-3, to=3-1]
    	\arrow["{R|_{\Sigma_t}\times R|_{\Sigma_t} \times \di R|_{\Sigma_t} \times \di R|_{\Sigma_t}}"{description}, from=1-3, to=3-5]
	    \arrow["{F_{R'}}"{description}, from=3-1, to=5-1]
	    \arrow["{\di R|_{\Sigma_t}}"{description}, from=3-3, to=5-1]
    	\arrow["{\di R|_{\Sigma_t}}"{description}, from=3-3, to=5-5]
	    \arrow["{F_R}"{description}, from=3-5, to=5-5]
    \end{tikzcd}\]
\end{axiom}

Trovare le espressioni per le forze più appropriate a modellare i problemi concreti in natura è il lavoro dei fisici: queste dipendono dalle caratteristiche specifiche dei sistemi considerati. Vediamo due esempi:

\begin{example}{Forza elettrostatica e forza gravitazionale}{coulomb e gravità}
    Due casi importanti di forze sono la forza \bemph{elettrostatica} che agisce su due punti con cariche elettriche $q, q' \in \R$ e posizioni $\bar x, \bar x' \in E^3$:
    \[ F_R(\bar x, \bar x', \bar v, \bar v') = k q q' \frac{ \bar x-\bar x' }{\| \bar x - \bar x' \|^3} \]
    E la forza \bemph{di gravità}, che dipende dalle masse invece:
    \[ F_R(\bar x, \bar x', \bar v, \bar v') = -G m m' \frac{ \bar x-\bar x' }{\| \bar x - \bar x' \|^3} \]
    Le costanti $k>0$ e $G>0$ sono dette rispettivamente \bemph{costante di Coulomb} e \bemph{costante di gravitazione universale}.\\
    Queste forze hanno una forma molto simile, la differenza cruciale è che la forza di gravità è sempre attrattiva, mentre quella elettrostatica è repulsiva tra cariche di segno uguale e attrattiva tra cariche di segno opposto.
\end{example}

Euristicamente, possiamo osservare che per non violare il primo principio ogni forza fisicamente sensata deve decrescere fino ad annullarsi al crescere della distanza, come le forze che abbiamo appena visto.\\
Nella nostra formulazione, il terzo principio è un'immediata conseguenza del secondo principio e di quello dell'impulso: infatti per due punti in moto rettilineo uniforme, sommando le due equazioni in \cref{ax:dinamica 2.1} otteniamo l'identità
\[ F_R(\gamma_Q|_R(t), \gamma_{Q'}|_R(t), v_Q|_R(t), v_{Q'}|_R(t)) = - F'_R(\gamma_{Q'}|_R(t), \gamma_{Q}|_R(t), v_{Q'}|_R(t), v_{Q}|_R(t))\]
Che viene generalizzata dal seguente principio, che permette di derivare anche la conservazione del momento angolare

\begin{axiom}{Principio forte di azione-reazione}{dinamica 2.2}
    Per ogni $t \in I$, se il vettore $\gamma_Q|_R(t) - \gamma_{Q'(t)}|_R(t)$, le forze 
    \[ F_R(\gamma_Q|_R(t), \gamma_{Q'}|_R(t), v_Q|_R(t), v_{Q'}|_R(t)) \quad\text{e}\quad F'_R(\gamma_{Q'}|_R(t), \gamma_{Q}|_R(t), v_{Q'}|_R(t), v_{Q}|_R(t)) \]
    Sono dirette lungo la direzione definita da tale vettore e hanno senso opposto.
\end{axiom}

Il prossimo principio ci permette di generalizzare quanto detto per i sistemi con due particelle a sistemi di un numero finito di particelle

\begin{axiom}{Principio di sovrapposizione delle forze}{sovrapposizione delle forze}
    Dato un sistema isolato di $N \in \N$ punti $\{Q_k\}_{k\le N}$, per ogni sistema di riferimento inerziale $R$ vale il secondo principio della dinamica in forma completa, ovvero per ogni $k\le N$ e per ogni $t \in I$:
    \[ m_{Q_k}a_{Q_k}|_R(t) = \underbrace{\sum_{i\neq k} F_{R,ki} (\gamma_{Q_k}|_R(t), \gamma_{Q_i}|_R(t), v_{Q_k}|_R(t), v_{Q_i}|_R(t)) }_{=: F_{R,k}(\gamma_{Q_1}|_R(t),\ ...\ ,\ \gamma_{Q_N}|_R(t),v_{Q_1}|_R(t),\ ...\ ,\  v_{Q_N}|_R(t)) } \]
    Dove $F_{R,ki}$ è la forza che agisce su $Q_k$ a causa di $Q_i$, ottenuta allontanando sufficientemente tutti gli altri punti
\end{axiom}

Mettendo l'equazione dataci da questo principio in coordinate, otteniamo un familiare sistema di equazioni differenziali del secondo ordine:
\[ \frac{\di^2 \bar x_k(t) }{\di t^2} = \frac{1}{m_k}  \bar F_k \left(\bar x_1(t),\ ...\ ,\ \bar x_N(t), \frac{\di \bar x_1(t)}{\di t},\ ...\ ,\  \frac{\di \bar x_N(t)}{\di t}\right) \]
Il \bemph{problema fondamentale della dinamica} è risolvere questo sistema date le condizioni iniziali e così prevedere il moto del sistema fisico.\\
È un problema difficile, ma il seguente teorema (la cui dimostrazione esula dagli scopi di questo corso) ci garantisce almeno il fatto che una soluzione esista e che la nostra teoria sia deterministica:

\begin{theorem}{Esistenza e unicità locale delle soluzioni del PDC}{esistenza e unicità locale 1}
    Siano $D,D' \subset \R^n$ aperti e $I \subset \R$ un intervallo aperto e si consideri l'equazione differenziale 
    \[ \frac{\di^2 \bar x(t)}{\di t^2} = F\left( t, \bar x(t), \frac{\di \bar x(t)}{\di t} \right) \]
    Con $F : I\times D, D'\to\R^n$ una funzione $\C^1$.\\
    Per ogni dato iniziale $(t_0, \bar x_0, \dot{\bar x}_0)$ e un intervallo aperto $I_0\subset I$ sufficientemente piccolo che contenga $t_0$, esiste ed è unica la funzione $\bar x:I_0\to\R^n$ di classe $\C^2$ che risolve l'equazione differenziale e soddisfa $\bar x(t_0) = \bar x_0$ e $\dot{ \bar x}(t_0) = \dot {\bar x}_0$\\
    Inoltre, ciascun dato iniziale determina un intervallo massimale in cui questa soluzione esiste ed è unica.
\end{theorem}

Facciamo qualche osservazione sull'invarianza della meccanica Newtoniana rispetto alle trasformazioni Galileiane, che sappiamoe essere gli automorfismi affini di $\V^4$.\\
Ricordiamo che in un sistema di riferimento inerziale $R$, le trasformazioni di Galileo prendono la forma:
\[ (t,\bar x)\mapsto (c+at, \bar b + \bar v t + A\bar x) \]
Per $c \in \R$, $a \in \{\pm 1\}$, $\bar b, \bar v \in \R^3$ e $A \in \O(3)$, opportunamente ristretti nel caso orientato.

\begin{axiom}{Principio di relatività Galileiana}{relatività Galileiana}
    Tutta la meccanica Newtoniana, in particolare le equazioni di \cref{ax:sovrapposizione delle forze} sono invarianti per le trasformazioni di Galileo.
\end{axiom}

Osserviamo che questo principio restringe la forma che possono avere le forze: infatti dato che una trasformazione di Galileo agisce su di una curva $t\mapsto\bar x(t)$ mandandola in $t' = \mapsto \bar x'(t') $, dove sostituende $t' = c + at$ vediamo $\bar x'(t') = \bar b + \bar v t + A\bar x(t)$ e dunque 
\[ \frac{\di \bar x'(t')}{\di t'} = a\bar v + a A \frac{\di \bar x(t)}{\di t} \qquad\text{e}\qquad \frac{\di^2 \bar x'(t')}{\di t'^2} = A\frac{\di^2 \bar x(t)}{\di t^2}\]
Per l'invarianza Galileiana delle leggi di Newton, otteniamo che una forza deve rispettare 
\[ \bar F\left( \bar x'(t'), \bar y'(t'), \frac{\di \bar x'(t')}{\di t'}.\frac{\di \bar y'(t')}{\di t'} \right) = A \bar F \left( \bar x(t), \bar y(t), \frac{\di \bar x(t)}{\di t},\frac{\di \bar y(t)}{\di t} \right)\]
per ogni trasformazione di Galileo. Questo implica che $\bar F$ deve essere della forma $\bar F(\bar x, \bar y, \bar u, \bar w) = \bar f(\bar x-\bar y, \bar u - \bar w)$ con $\bar f : \R^3\times \R^3 \to \R^3$ che rispetta le seguenti proprietà:\begin{itemize}
    \item $\bar f(A\bar x , A\bar u) = A\bar f(\bar x, \bar u)$ per ogni $A \in \O(3)$.
    \item $\bar f(\bar x, a\bar u) = \bar f(\bar x, \bar u)$ per $a =\pm 1$
\end{itemize}

\subsection{Quando Newton non funziona}

I principi appena formulati non contemplano numerosi casi di interesse fisico, di cui adesso presentiamo una piccola selezione

\newcommand\eff{\text{eff}}
\subsubsection{Moto assegnato per un sottosistema e forze dipendenti dal tempo}

Molto frequentemente, quasi sempre in effetti, nelle situazioni pratiche capita che in un sistema di punti $\{Q_1,\ ...\ ,\ Q_N\}$ il moto di un sottosistema $\{Q_{n+1},\ ...\ ,\ Q_N\}$ sia fissato e si voglia trovare di conseguenza il moto dei punti $\{Q_1,\ ...\ ,\ Q_n\}$ rimanenti; è il caso ad esempio del moto della Terra fissato quello del Sole.\\
Inserendo le linee di universo prescritte di $\{Q_{n+1},\ ...\ ,\  Q_N\}$ nelle equazioni di Newton si trova una forza "effettiva" che può dipendere esplicitamente dal tempo:
\[F^\eff_k(t, \gamma_{Q_1}(t),\ ...\ ,\  \gamma_{Q_n}(t), v_{Q_1}(t),\ ...\ ,\ v_{Q_n}(t)) := F_k(\gamma_{Q_1}(t),\ ...\ ,\  \gamma_{Q_N}(t),  v_{Q_1}(t),\ ...\ ,\ v_{Q_N}(t) )\]
Le linee di universo dei punti rimanenti sono quindi determinate dal sistema 
\[ m_{Q_k} a_{Q_k}(t) = F^\eff_k(t, \gamma_{Q_1}(t),\ ...\ ,\  \gamma_{Q_n}(t), v_{Q_1}(t),\ ...\ ,\ v_{Q_n}(t)) \]
che comunque soddisfa le ipotesi del teorema di esistenza e unicità delle soluzioni, dunque possiamo estendere il determinismo a situazioni più generali.

\newcommand\cent{\text{(cent)}}
\subsubsection{Sistemi non inerziali}

A volte conviene usare un sistema di riferimento non inerziale $R'$, ad esempio un sistema solidale alla superficie della Terra, ma questo modifica un po' la forma delle equazioni di Newton. Dalle leggi di trasformazione ottengo le nuove equazioni 
\[ m_{Q_k}a_{Q_k}|_{R'}(t) + m_{Q_k}a_{Qk}^\tv(t) + m_{Q_k}a_{Q_k}^\cor (t) = F_{R',k}\left( t, \gamma_{Q_1}|_{R'}(t),..,\gamma_{Q_N}|_{R'}(t), v_{Q_1}|_{R'}(t),\ ...\ ,\  v_{Q_n}|_{R'}(t)  \right) \]
ricordando che $v|_R(t) = v|_{R'}(t)+v^\tv(t)$ e $a|_R(t) = a|_{R'}+a^\tv(t)+a^\cor(t)$: spesso i due termini "aggiuntivi" a sinistra vengono portati a destra e interpretati come le cosiddette forze fittizie, dovute al sistema non inerziale:
\[ m_{Q_k}a_{Q_k}|_{R'} = \underbrace{F^\cent}_{=: -m_{Q_k}a_{Q_k}^\tv} + \underbrace{F^\cor}_{=: -m_{Q_k}a_{Q_k}^\cor} + F_{R', k} \]
Dato che il lato destro comunque dipende da $t$, $\gamma_{Q_i}$ e $v_{Q_i}$ si applica comunque il teorema di esistenza e unicità dandoci il determinismo.

\subsubsection{Vincoli geometrici}

Nelle applicazioni reali, spesso capita che i punti del sistema $\{Q_1,\ ...\ ,\ Q_N\}$ debbano muoversi rispettando certi vincoli geometrici: magari la posizione di $Q_3$ deve trovarsi su una qualche curva perchè scorre su un binario, o $Q_1$ e $Q_5$ devono trovarsi a distanza fissata perchè uniti da un'asta, o tutti i punti devono giacere sulla superficie di un tavolo, chi più ne ha più ne metta.\\
Si nota immediatamente che le equazioni di Newton viste sopra non sono adeguate a descrivere questi sistemi vincolati, perchè la dinamica indotta non presenta vincoli geometrici: per affrontare questo problema si introducono le cosiddette \bemph{reazioni vincolari}, delle "forze" $\phi_k$ che modificano le equazioni del moto nella forma
\[ m_{Q_k} a_{Q_k}|_R(t) = F_{R,k}\left( \gamma_{Q_1}|_{R}(t),..,\gamma_{Q_N}|_{R}(t), v_{Q_1}|_{R}(t),\ ...\ ,\  v_{Q_n}|_{R}(t)  \right) + \phi_k \]
Ma queste reazioni vincolari non sono note a priori, appaiono come incognite nel sistema e questo in generale ci porta fuori dalle ipotesi sotto cui ci è garantito il determinismo dagli strumenti sviluppati fino a questo punto.\\ 
Per trattare questi sistemi introdurremo il formalismo della meccanica Lagrangiana, che svilupperemo nella prossima parte del corso.

\chapter{Meccanica Lagrangiana}

Il formalismo della meccanica Lagrangiana è una formulazione più elegante e soprattutto estremamente più potente della meccanica classica rispetto alla meccanica Newtoniana.\\
Alcuni dei principali vantaggi sono:\begin{itemize}
    \item Questa formulazione si applica anche a certe classi di sistemi vincolati, dove i vincoli sono detti \bemph{ideali}, estendendo consistentemente la meccanica Newtoniana.
    \item Il legame tra le simmetrie di un sistema e le sue quantità conservate datoci dal teorema di Noether, che vedremo più avanti, viene espresso in modo estremamente elegante e naturale.
\end{itemize}

\section{Geometria dei vincoli olonomi e ideali}

In questa sezione ci occuperemo dei sistemi di $N\ge 1$ punti materiali $\{Q_1,\ ...\ ,\ Q_N\}$. Dato che questi $N$ punti sono descritti da $N$ linee di universo $\gamma_k : I \to \V^4$ con $T\circ \gamma_k = \id_{I}$, per tenere traccia di tutti questi punti contemporaneamente introduciamo un'estensione del nostro spaziotempo.

\begin{definition}{Spaziotempo delle configurazioni}{spaziotempo delle configurazioni}
    Sia $T:\V^4\to E^1$ uno spaziotempo e sia $N\ge 1$.\\
    Il corrispondete \bemph{spaziotempo delle configurazioni} per $N$ punti materiali è dato dal prodotto fibrato 
    \[ \V^{3N+1} = \underbrace{\V^4 \times_{\E_1}...\times_{\E_1} \V^4 }_{N\text{ volte}} = \left\{ (p_1,\ ...\ ,\ p_N) \in (\V^4)^N : \forall i, j, T(p_i) = T(p_j) \right\} \]
    Con la sommersione suriettiva ovvia $T:\V^{3N+1}\to E^1$ data da $T(p_1,\ ...\ ,\ p_N) = T(p_i)$ per qualsiasi $i$.\\
    Alternativamente, $T:\V^{3N+1}\to E^1$ è definito per proprietà universale come limite del seguente diagramma nella categoria dei fibrati differenziabili 
    \[\begin{tikzcd}
    	& {\V^{3N+1}} & \\
    	{\V^4} && {\V^4} \\
	    & {E^1}
    	\arrow[dashed, from=1-2, to=2-1]
	    \arrow[dashed, from=1-2, to=2-3]
    	\arrow["\lrcorner"{anchor=center, pos=0.125, rotate=-45}, draw=none, from=1-2, to=3-2]
	    \arrow["{...}"{description}, draw=none, from=2-1, to=2-3]
    	\arrow["T"{description}, from=2-1, to=3-2]
	    \arrow["T"{description}, from=2-3, to=3-2]
    \end{tikzcd}\]
\end{definition}

Osserviamo che le fibre di $\V^{3N+1}$ al tempo $t\in E^1$ sono date da $\Sigma_t^N = \Sigma_t\times...\times\Sigma_t$, che dunque eredita la struttura di spazio euclideo sopra lo spazio vettoriale sottostante $V_t^N = V_t\times...\times V_t$ con l'evidente prodotto scalare.\\
Trattare un sistema di $N$ linee di universo $\gamma_{k}:I\to\V^4$ è assolutamente identico al trattare una singola linea di universo $\Gamma :I\to\V^{3N+1}$ data da $\Gamma(t)=(\gamma_1(t),\ ...\ ,\ \gamma_N(t))$, infatti come visto sopra $T\circ \Gamma = \id_{I}$ e viceversa proiettando su ciascuna componente torniamo a $\gamma_k = \pi_k\circ \Gamma: I \to \V^4$.\\
Anche un sistema di riferimento $R:\V^4\to E^1\times E^3$ può essere facilmente esteso e illuminare un attimo il senso di quel $3N+1$ all'esponente, come si vede in questo diagramma 
\[\begin{tikzcd}
	{\V^4\times_{\E_1}...\times_{\E_1}\V^4} &&& {(E^1\times E^3)\times_{E^1}...\times_{E^1}(E^1\times E^3)} \\
	{\V^{3N+1}} &&& {E^1\times E^{3N}} \\
	\\
	{E^1} &&& {E^1}
	\arrow["{(R,\ ...\ ,\  R)}"{description}, from=1-1, to=1-4]
	\arrow["\cong"{marking, allow upside down}, draw=none, from=1-4, to=2-4]
	\arrow["{=}"{marking, allow upside down}, draw=none, from=2-1, to=1-1]
	\arrow["T"{description}, from=2-1, to=4-1]
	\arrow["{\pi_{E^1}}"{description}, from=2-4, to=4-4]
	\arrow["{\id_{E^1}}"{description}, from=4-1, to=4-4]
\end{tikzcd}\] 
E lo stesso vale evidentemente per qualsiasi scelta di coordinate standard che ci porterà in $\R^1\times\R^{3N}$.

\begin{definition}{Vincoli olonomi}{vincoli olonomi}
    Un sistema di $0\le C\le N$ \bemph{vincoli olonomi} per un sistema di di $N\ge 1$ punti materiali è dato da una famiglia di $C$ funzioni lisce $f_j : \V^{3N+1}\to\R$ che soddisfano le seguenti proprietà:\begin{itemize}
        \item Per ogni $t \in E^1$ vale 
        \[ \Sigma_t^N \cap \bigcap_j f_j^{-1}(\{0\}) \neq \varnothing\]
        \item Per ogni $t \in E^1$ e qualsiasi scelta di coordinate standard adattate $(t,\bar{\bar x}) := (t,x^1,\ ...\ ,\ x^{3N})$, la matrice Jacobiana 
        \[ \left[ \frac{\partial f_j}{\partial x^i} \bigg|_{(t,\bar{\bar x})} \right]_{1\le i\le 3N,1\le j\le C} : \R^{3N}\to \R^C \]
        Ha rango massimo per ogni $(t,\bar{\bar x})$ tale che $f_j((t,\bar{\bar x})) = 0$ per ogni $j\le C$
    \end{itemize}
    A questo punto è detto \bemph{spaziotempo delle configurazioni} di un sistema di $N$ punti materiali sottoposto a $C$ vincoli olonomi $\{f_1,\ ...\ ,\ f_C\}$ il fibrato 
    \[ T: \V^{n+1}:= \bigcap_j f_j^{-1}(\{0\}) \subset \V^{3N+1} \to E^1 \]
    Dove $n := 3N - C$ è detto il numero di \bemph{gradi di libertà} del sistema.\\
    Lo spazio $\Q_t := \V^{n+1}\cap \Sigma_t^N$ è detto \bemph{spazio delle configurazioni} del sistema al tempo $t$.
\end{definition}

Ok, che diavolo abbiamo appena letto? Le proprietà che abbiamo richiesto sui vincoli sono proprietà tecniche che ci permettono di applicare un potente teorema di geometria differenziale, il \bemph{teorema dei valori regolari}\footnote{Reperibile in \cite{Moretti2024} per interessati e masochisti} che ci conduce al seguente risultato.

\begin{proposition}{Utile sulle varietà}{utile sulle varietà}
    Nel contesto della definizione precedente, valgono i seguenti fatti:\begin{enumerate}
        \item $\V^{n+1}\subset \V^{3N+1}$ è una sottovarietà \textit{embedded}\footnote{Possibili traduzioni sono immersa, inclusa o incastonata, ma in genere si usa l'aggettivo in inglese e ci si mette il cuore in pace.} di dimensione $n+1$ e la restrizione $T:\V^{n+1}\to E^1$ è una sommersione suriettiva.
        \item Per ogni $t \in E^1$, $\Q_t = \V^{n+1}\cap\sigma_t^N \subset \Sigma_t^N$ è una sottovarietà embedded di dimensione $n$.
        \item Per ogni punto $p \in \V^{3N+1}$ possiamo scegliere in un intorno di $p$ una carta locale $(U\subset \V^{3N+1},\phi:U\to \R^{3N+1})$ tale che \begin{itemize}
            \item \[ \forall p' \in U, \phi(p') := (T(p'), \phi^1(p'),\ ...\ ,\ \phi^{3N}(p')) \]
            \item \[ \forall t \in T(U), \forall p' \in U\cap Q_t, \phi(p') = (t, \phi^1(p'),\ ...\ ,\ \phi^n(p'),0,\ ...\ ,\ 0) \]
        \end{itemize}
    \end{enumerate}
    \proof 
    La seconda proprietà segue direttamente dalle richieste sui vincoli olonomi e dal teorema dei valori regolari; per verificare la prima proprietà dobbiamo estendere la matrice Jacobiana con una derivata temporale, che faremo identificando $x^0:= t$:
    \[\left[ \frac{\partial f_j}{\partial x^i} \bigg|_{(t,\bar{\bar x})} \right]_{0\le i\le 3N,1\le j\le C} : \R^{3N+1}\to \R^C\]
    questa matrice ha ovviamente rango massimo, perchè restringendoci a vettori della forma $(0,\bar{\bar{x}})$ otteniamo la matrice Jacobiana non estesa, che ha rango massimo per definizione.\\
    Rimane solo la terza, che segue dal teorema dei valori regolari.
\end{proposition}

Vediamo un esempio

\begin{example}{Un punto su una sfera}{un punto su una sfera}
    Consideriamo un sistema di un ($N=1$) punto materiale e un ($C=1$) vincolo: scegliendo un qualsiasi sistema di riferimento $R$, abbiamo $\V^{3+1}\cong E^1\times E^3$ possiamo considerare la mappa liscia $f:E^1\times E^3\to\R$ definita da $(t, q)\mapsto \| q- O(t) \|^2 - r_0^2(t)$, dove $O:E^1\to E^3$ e $r_0:E^1\to\R_{>0}$ sono mappe lisce.\\
    Per ogni $t$ fissato, i punti di $f(t,q)=0$ definiscono una sfera $\S^2_t$ di centro $O(t)$ e raggio $r_0(t)$. Passando in coordinate abbiamo:
    \[ f(t,\bar x) = (x^1 - O^1(t))^2 + (x^2 -O^2(t))^2 + (x^3 -O^3(t))^2 - r_0^2(t) \]
    Chiaramente il primo punto della definizione di vincolo olonomo è soddisfatto dal fatto che la controimmagine di $0$ è come abbiamo detto una sfera; per la seconda proprietà dobbiamo calcolare la matrice Jacobiana:
    \[ \left[ \frac{\partial f}{\partial x^1}, \frac{\partial f}{\partial x^2}, \frac{\partial f}{\partial x^3} \right] = \left[ 2(x^1 -O^1(t)), 2(x^2 -O^2(t)), 2(x^3 -O^3(t)) \right] \]
    Che ha rango massimo fintanto che almeno una delle sue componenti non è nulla, cosa che si verifica sempre negli zeri di $f$, dunque questo è un vincolo olonomo.
\end{example}

% \begin{example}{asdasd}{asfsaf}
%     Consideriamo ora $N=2$ punti materiali, dunque con un riferimento $R$ ci poniamo in $\V^{6+1}\iso E^1\times E^3\times E^3$ e le mappe lisce $f_{1,2}:E^1\times E^3\times E^3$ definite da 
%     \[ f_1(t,q_1,q_2) = \| q_1-O(t) \|^2 - r_0^2(t)\]
%     \[ f_2(t,q_1,q_2) = \|q_1-q_2\|^2 - D_0^2(t) \]
%     TBD
% \end{example} 

In preprarazione alla discussione sui vincoli ideali, definiamo un importante concetto di geometria differenziale:

\begin{definition}{Spazio tangente di $\Q_t$}{spazio tangente}
    Consideriamo, per $t\in\Q^1$, la sottovarietà embedded $\Q_t = \V^{n+1}\cap\Sigma^N_t \subset \Sigma^N_t$.\\
    Scegliendo coordinate locali su $\Q_t$, ovvero una carta locale $(U,\psi)$ tale che $U\subset\Q_t$ e $\psi:U\to\R^n$ otteniamo una parametrizzazione locale $\psi^{-1}:\psi(U)\to\Sigma^N_t$ di $\Q_t$ data da $\psi^{-1}(q^1,\ ...\ ,\ q^n)$.\\
    Definiamo lo \bemph{spazio tangente} di $\Q_t$ in $q \in U$ come il sottospazio vettoriale 
    \[ T_q\Q_t := \operatorname{span}_{\R}\left( \frac{\partial \psi^{-1}(q^1,\ ...\ ,\ q^n)}{\partial q^1}\bigg|_{\psi(q)},\ ...\ ,\ \frac{\partial \psi^{-1}(q^1,\ ...\ ,\ q^n)}{\partial q^n}\bigg|_{\psi(q)}  \right) \]
    di $V_t^N$, spazio sottostante a $\Sigma_t^N$.
\end{definition}

Dimostriamo alcuni utili risultati su $T_q\Q_t$:

\begin{proposition}{Essenziale unicità dello spazio tangente}{essenziale unicità dello spazio tangente}
    Lo spazio tangente $T_q\Q_t$ non dipende dalla scelta di coordinate e come spazio vettoriale ha la stessa dimensione $n$ di $\Q_t$ come varietà.
    \proof 
    Senza esplicitare tutti i calcoli, date due carte locali $(U,\psi)$ e $(U',\psi')$ si nota che la funzione di trasferimento $\psi'\circ\psi^{-1}$; calcolando con la regola della catena le derivate di $\psi^{-1} = \psi'^{-1}\circ \psi'\circ \psi^{-1}$ otteniamo 
    \[ \frac{\partial \psi^{-1}(q^1,\ ...\ ,\  q^n)}{\partial q^i}\bigg|_{\psi(q)} = \sum_j \frac{\partial \psi'^{-1}(q'^1,\ ...\ ,\  q'^n)}{\partial q'^j}\bigg|_{\psi(q)}\frac{\partial q'^j}{\partial q^i}\bigg|_{\psi(q)} \]
    Poichè la Jacobiana di ogni diffeomorfismo è invertibile, ne segue che gli spazi tangenti coincidono.\\
    Per dimostrare che la dimensione è la stessa di $\Q_t$, usiamo una carta del tipo del terzo punto di \cref{prop:utile sulle varietà}: in questo caso è evidente che i generatori di $T_q\Q_t$ sono indipendenti, dato che usando coordinate locali su $\Sigma_t^N$ si ottiene $\psi^{-1}(q^1,\ ...\ ,\ q^n) = (q^1,\ ...\ ,\ q^n,0,\ ...\ ,\ 0)$ e derivando rispetto a $q^i$ si osserva l'indipendenza lineare.
    \qed 
\end{proposition}

Finalmente possiamo definire i vincoli ideali:

\begin{definition}{Vincoli ideali}{vincoli ideali}
    Sia $\{f_1,\ ...\ ,\  f_C\}$ un sistema di $0\le C\le 3N$ vincoli olonomi per un sistema di $N\ge 1$ punti materiali.\\
    Questi vincoli sono detti \bemph{ideali} se per ogni linea di universo $\Gamma:I\to\V^{n+1}$ e ogni $t\in I$, la reazione vincolare $\phi = (\phi_1,\ ...\ ,\ \phi_N) \in V_t^N$ è normale a $\Q_t$, ovvero $\phi \cdot v = 0$ per ogni vettore tangente $v \in T_{\Gamma(t)}\Q_t$
\end{definition}

Osserviamo che questa definizione ha una forte connotazione fisica, infatti questi vincoli sono stati introdotti nel contesto della meccanica Newtoniana: in particolare, non si può dedurre dalla definizione matematica di vincoli olonomi che questi debbano essere o non essere anche ideali.\\
Si deve pensare alla richiesta di idealità come a una richiesta aggiuntiva sul comportamento geometrico di un vincolo che lo rende adatto per applicare l'apparato della meccanica Lagrangiana.

\newpage
\newcommand\T{\mc{T}}
\section{Energia cinetica}

Per trattare i sistemi di punti materiali con vincoli ideali è spesso necessario calcolare le componenti tangenziali ai vincoli dell'accelerazione, ma per farlo è necessario introdurre il concetto di energia cinetica.\\
L'obiettivo di questa parte sarà definire l'energia cinetica per sistemi di $N\ge 1$ punti materiali $\{Q_1,\ ...\ ,\ Q_N\}$ sottoposti a $C\le N$ vincoli olonomi $\{f_1,\ ...\ ,\ f_C\}$ e illustrare la geometria sottesa alle coordinate puntate $\dot q$. Per farlo dobbiamo estendere opportunamente i concetti di velocià e accelerazione in $\V^4$ a $\V^{3N+1}$\\
Sia $\Gamma:I\to \V^{n+1}$ una linea di universo nello spazio delle configurazioni vincolato: dato che questo è una sottovarietà embedded di $\V^{3N+1}$, possiamo leggerla come linea di universo in quest'ultimo e definire come sopra velocità e accelerazione rispetto a un riferimento $R$: esplicitamente
\[ v|_R(t) := \left((\di R|_{\Sigma_t}^{\times N})^{-1}\circ \frac{\di}{\di t} \circ R|_{\Sigma_t}^{\times N} \right)(\Gamma(t)) = (v_{Q_1}|_R(t),\ ...\ ,\  v_{Q_N}|_R(t)) \in V_t^\N \]
\[ a|_R(t) := \left((\di R|_{\Sigma_t}^{\times N})^{-1}\circ \frac{\di^2}{\di t^2} \circ R|_{\Sigma_t}^{\times N} \right)(\Gamma(t)) = (a_{Q_1}|_R(t),\ ...\ ,\  a_{Q_N}|_R(t)) \in V_t^\N \]
Ora possiamo definire l'energia cinetica.

\begin{definition}{Energia cinetica}{energia cinetica}
    Sia $\Gamma:I\to\V^{n+1}\subset\V^{3N+1}$ una linea cinetica nello spaziotempo delle configurazioni di un sistema con vincoli olonomi e sia $R$ un riferimento.\\
    La sua \bemph{energia cinetica} rispetto a $R$ è definita da:
    \[ \T_R(t) := \sum_i \frac{m_{Q_i}}{2}\| v_{Q_i}|_R(t) \|^2 =  \sum_i \frac{m_{Q_i}}{2}(v_{Q_i}|_R(t) \cdot v_{Q_i}|_R(t)) \]
    Dove il prodotto scalare è da intendersi come il prodotto dello spazio euclideo $\Sigma_t$
\end{definition}

Sarà opportuno ottenere una formula per l'energia cinetica in coordinate.\\
Data una carta $(U,\psi)$, questa si dice \bemph{adattata alla sommersione suriettiva} $T:\V^{n+1}\to E^1$ se è della forma $\psi(q) = (T(q),q^1(q),\ ...\ ,\ q^n(q))$; le coordinate locali $(t,q^1,\ ...\ ,\ q^n)$ sono dette \bemph{coordinate locali naturali} su $\V^{n+1}$.\\
Notiamo che la restrizione $\phi|_{\Sigma_t^N}$ definisce una carta locale su $\Q_t$ per $t$ fissato; per $I$ sufficientemente piccoli, la linea di universo $\Gamma : I\to\V^{n+1}$ può essere fattorizzata attraverso una carta locale adattata $\psi$ in $\psi\circ\Gamma : I \to \psi(U)\subset \R^{n+1}$ con $\psi\Gamma(t) = (t,q^1(t),\ ...\ ,\ q^n(t))$; allo stesso modo una mappa liscia $\gamma:I\to\psi(U)\subset\R^{n+1}$ della forma $t\mapsto(t,q^1(t),\ ...\ ,\ q^n(t))$ definisce una linea di universo attraverso la composizione $\psi^{-1}\circ\gamma$.

\begin{lemma}{}{}
    Sia $\Gamma:I\to\V^{n+1}$ una linea di universo con $I$ sufficientemente piccolo da permettere una rappresentazione in coordinate locali naturali $\Gamma(t) = \psi^{-1}(t,q^1(t),\ ...\ ,\ q^n(t))$. Allora la sua velocità può essere scritta come 
    \[ v|_R(t) = \frac{\partial }{\partial t}\bigg|_R \psi^{-1}(t, q^1(t),\ ...\ ,\ q^n(t)) +\sum_k \frac{\partial \psi|_{\Sigma_t^N}^{-1}(q^1(t),\ ...\ ,\ q^n(t))}{\partial q^k}\frac{\di q^k(t)}{\di t} \]
    Dove il primo termine è definito da 
    \[ \frac{\partial }{\partial t}\bigg|_R = (\di R|_{\Sigma_t}^{\times N})^{-1}\circ \frac{\partial}{\partial t}\circ R|_{\Sigma_t}^{\times N} \]
    \proof 
    Usiamo la formula della velocità che abbiamo definito sopra:
    \begin{align*}
        v|_R(t) & = \left((\di R|_{\Sigma_t}^{\times N})^{-1}\circ \frac{\di}{\di t} \circ R|_{\Sigma_t}^{\times N} \right)(\psi^{-1}(t,q^1(t),\ ...\ ,\ q^n(t))) \\ & = (\di R|_{\Sigma_t}^{\times N})^{-1}\circ\left[ \frac{\partial R|_{\Sigma_t}^{\times N}(\psi^{-1}(...))}{\partial t} + \sum_k \underbrace{\frac{\partial R|_{\Sigma_t}^{\times N}(\psi^{-1}(...))}{\partial q^k}}_{= \di R|_{\Sigma_t}^{\times N} \frac{\partial \psi|_{\Sigma_t^N}^{-1}(...)}{\partial q^k}}\frac{\di q^k(t)}{\di t} \right] \\ & = \left((\di R|_{\Sigma_t}^{\times N})^{-1}\circ \frac{\partial}{\partial t}\circ R|_{\Sigma_t}^{\times N}\right)(\psi^{-1}(...)) + \sum_k \frac{\partial \psi|_{\Sigma_t^N}^{-1}(q^1(t),\ ...\ ,\ q^n(t))}{\partial q^k}\frac{\di q^k(t)}{\di t}
    \end{align*}
    \qed 
\end{lemma}

Notiamo che le derivate parziali rispetto ai vari $q^k$ di $\psi|_{\Sigma_t^N}^{-1}(...)$ sono vettori tangenti a $\Q_t$ che hanno $N$ componenti in $V_t$, ciascuna accessibile tramite la proiezione all'$i$-esimo fattore $\pi_i$.\\
Inserendo le velocità in coordinate nella formula dell'energia cinetica, otteniamo 
\begin{align*}
    \T_R(t) & = \sum_{k,l} A_{k,l}(t, q^1(t),\ ...\ ,\  q^n(t))\frac{\di q^k(t)}{\di t}\frac{\di q^l(t)}{\di t} \\ & +  \sum_k B_k(t, q^1(t),\ ...\ ,\  q^n(t)) \frac{\di q^k(t)}{\di t} \\ & + C(t, q^1(t),\ ...\ ,\  q^n(t))
\end{align*}
con 
\[ A_{k,l}(t, q^1(t),\ ...\ ,\  q^n(t)) = \sum_i \frac{m_{Q_i}}{2}  \left( \pi_i\left( \frac{\partial \psi|_{\Sigma_t^N}^{-1}(t,q^1(t),\ ...\ ,\ q^n(t)) }{\partial q^k} \right) \cdot \pi_i\left( \frac{\partial \psi|_{\Sigma_t^N}^{-1}(t,q^1(t),\ ...\ ,\ q^n(t)) }{\partial q^l} \right) \right) \]
\[ B_k(t, q^1(t),\ ...\ ,\  q^n(t)) = 2\sum_i \frac{m_{Q_i}}{2}  \left( \pi_i\left( \frac{\partial \psi|_{\Sigma_t^N}^{-1}(t,q^1(t),\ ...\ ,\ q^n(t)) }{\partial q^k} \right) \cdot \pi_i\left( \frac{\partial }{\partial t}\bigg|_R\psi^{-1}(t,q^1(t),\ ...\ ,\ q^n(t))  \right) \right) \]
\[ C(t, q^1(t),\ ...\ ,\  q^n(t)) = \sum_i \frac{m_{Q_i}}{2}  \left( \pi_i\left( \frac{\partial }{\partial t}\bigg|_R\psi^{-1}(t,q^1(t),\ ...\ ,\ q^n(t)) \right) \cdot \pi_i\left( \frac{\partial }{\partial t}\bigg|_R\psi^{-1}(t,q^1(t),\ ...\ ,\ q^n(t))  \right) \right) \]
Questa mostruosità ci permette di pensare all'energia cinetica come a una funzione $\T(t,q^1,\ ...\ ,\ q^n,\dot q^1,\ ...\ ,\ \dot q^n)$ in $2n+1$ variabili in $\psi(U)\times\R^n$, dove le $\dot q^i$ rappresentano le derivate temporali. Ma che oggetto geometrico sottosta a queste nuove coordinate?

\subsection{Spaziotempo degli atti di moto}

Ci sono diversi metodi avanzati e molto eleganti per costruire i fibrati di jet, ma noi purtroppo non li possediamo, quindi dobbiamo usare il \textit{metodo della bestia} (sic)\footnote{Moretti 2023}, più elementare ma anche più sgraziato.

\subsubsection{Costruzione di una varietà dalla scelta di un atlante}

Sia $\M$ una varietà e sia $\A := \{ (U_i,\psi_i) \}_i$ un atlante non necessariamente massimale. Lo spazio quoziente $\bigsqcup_i U_i/\sim$ ottenuto dall'unione disgiunta degli aperti modulo la relazione di equivalenza $[i,p]\sim[j, q]$ se e solo se $p=q$ è omeomorfo a $\M$ tramite la mappa $[i,p]\mapsto p$.\\
Usando gli omeomorfismi delle carte locali mandando $[i,p]\mapsto [i,\psi_i(p)]$, possiamo leggerlo come $ \bigsqcup_i \psi_i(U_i)/\sim $, dove $[i,x]\sim[j,y]$ se e solo se $\psi_i^{-1}(x) = \psi_j^{-1}(y)$, ovvero $x = \psi_i\psi_j^{-1}(y)$ scrivendolo in termini delle funzioni di trasferimento.\\
Trasferendo la struttura differenziabile di $\M$ lungo l'omeomorfismo $\bigsqcup_i \psi_i(U_i)/\sim \iso \M$, otteniamo una descrizione diffeomorfa di $\M$ in termini di carte locali incollate tra loro mediante le loro funzioni di trasferimento.

\begin{definition}{Spaziotempo degli atti di moto}{spaziotempo degli atti di moto}
    Sia $T:\V^{n+1}\to E^1$ uno spaziotempo delle configurazioni per un sistema di vincoli olonomi e sia $\A = \{(U_i,\psi_i)\}_i$ il suo atlante dato dalle carte locali adattate alla sommersione suriettiva $T$.\\
    Lo \bemph{spaziotempo degli atti di moto} per il sistema $A(\V^{n+1})$, o \bemph{fibrato di jet}, è definito applicando la costruzione precedente alle carte di coordinate locali $\phi_i(U_i)\times\R^n$ e alle mappe di transizione $\psi_j(U_i\cap U_j)\times\R^n\iso \psi_i(U_i\cap U_j)\times \R^n$ date da
    \[ (t,\bar q ,\bar{\dot q}) \mapsto \left( t, (\psi_i\psi_j^{-1})|_{\Sigma_t^N}(\bar q), \frac{\partial (\psi_i\psi_j^{-1})(t, \bar q) }{\partial t} + \sum_k \frac{\partial (\psi_i\psi_j^{-1})|_{\Sigma_t^N}(\bar q) }{\partial q^k} \dot q^k \right) \] 
\end{definition}

Notiamo che la regola di trasformazione delle coordinate puntate $\dot{\bar q}$ codifica precisamente il modo in cui la derivata temporale di una curva liscia $\bar q(t)$ si trasforma attraverso un cambiamento di coordinate secondo la regola della catena.\\
Abbiamo un'evidente sommersione suriettiva $T:A(\V^{n+1})\to \V^{n+1}\to E^1$ che manda $(t,\bar q ,\bar{\dot q})\mapsto(t,\bar q )\mapsto t$. Allo stesso modo possiamo estendere una linea di universo $\Gamma:I\to\V^{n+1}$ a una linea $\hat\Gamma:I\to A(\V^{n+1})$ aggiungendo le derivate delle componenti:
\[ \hat\Gamma(t):= \left(t, q^1(t),\ ...\ ,\ q^n(t), \frac{\di q^1(t)}{\di t},\ ...\ ,\  \frac{\di q^n(t)}{\di t}\right) \]
Notiamo anche che la velocità può essere considerata una funzione $A(\V^{n+1})\to V_t^N$ ponendo
\[ v|_R(t,\bar q,\dot{\bar q}) = \frac{\partial }{\partial t}\bigg|_R \psi^{-1}(t,\bar q,\dot{\bar q}) +\sum_k \frac{\partial \psi|_{\Sigma_t^N}^{-1}(t,\bar q,\dot{\bar q})}{\partial q^k}\dot q^k \]
Possiamo fare un ragionamento analogo per l'energia cinetica dunque 
\[ \T(t,\bar q,\dot{\bar q}) = \sum_i \frac{m_{Q_i}}{2} \| v_{Q_i}|_R(t,\bar q,\dot{\bar q}) \|^2 \]
E questo ci permette di recuperare le espressioni per la velocità ed energia cinetica di $\Gamma$ attraverso la sua estensione $\hat\Gamma$:
\[ v_R(t) = [v_R]_{\hat\Gamma(t)} := v_R(\hat\Gamma(t)) \]
\[ \T_R(t) = [\T_R]_{\hat\Gamma(t)} := \T_R(\hat\Gamma(t)) \]

\begin{proposition}{Pre-Eulero--Lagrange}{premessa ad Eulero--Lagrange}
    Sia $\Gamma:I\to\V^{n+1}$ una linea di universo con $I$ abbastanza piccolo da ammettere una rappresentazione in coordinate locali naturali.\\
    Allora
    \[ \frac{\di}{\di t}\left[ \frac{\partial \T_R}{\partial \dot q^k} \right]_{\hat\Gamma(t)} - \left[ \frac{\partial \T_R}{\partial q^k} \right]_{\hat\Gamma(t)} = \sum_i m_{Q_i} \left( a_{Q_i}|_R(t)\cdot \pi_i\left( \frac{\partial \psi|_{\Sigma_t^N}^{-1}(t,\bar q,\dot{\bar q})}{\partial q^k} \right) \right) \]
    \proof 
    Omettiamo gli argomenti delle funzioni per alleggerire la notazione\footnote{Potreste chiederci perchè non l'abbiamo fatto prima... bella domanda, diremo che dovevamo tutti interiorizzare quali fossero questi argomenti.} e sviluppiamo il lato sinistro: 
    \begin{align*}
        \frac{\di}{\di t}\left[ \frac{\partial \T_R}{\partial \dot q^k} \right]_{\hat\Gamma(t)} - \left[ \frac{\partial \T_R}{\partial q^k} \right]_{\hat\Gamma(t)} & = \sum_i \frac{m_{Q_i}}{2} \left( \frac{\di}{\di t}\left[ \frac{\partial (v_{Q_i}|_R\cdot v_{Q_i}|_R) }{\partial \dot q^k} \right]_{\hat\Gamma(t)} - \left[ \frac{\partial (v_{Q_i}|_R\cdot v_{Q_i}|_R) }{\partial q^k} \right]_{\hat\Gamma(t)} \right) \\ 
        & = \sum_i m_{Q_i} \left( \frac{\di}{\di t} \left[ v_{Q_i}|_R\cdot \frac{\partial v_{Q_i}|_R }{\partial\dot q^k} \right]_{\hat\Gamma(t)} - \left[ v_{Q_i}|_R\cdot \frac{\partial v_{Q_i}|_R }{\partial q^k} \right]_{\hat\Gamma(t)} \right) \\ 
        & = \sum_i m_{Q_i} \left( \left(a_{Q_i}|_R \cdot \frac{\partial v_{Q_i}|_R}{\partial \dot q^k} \right) + \left( v_{Q_i}|_R \cdot \frac{\di}{\di t}\bigg|_R \frac{\partial v_{Q_i}|_R}{\partial \dot q^k} \right) - \left(v_{Q_i}|_R \cdot \frac{\partial v_{Q_i}|_R}{\partial q^k} \right) \right)\\
        & = \sum_i m_{Q_i}\left( \left(a_{Q_i}|_R \cdot \underbrace{\frac{\partial v_{Q_i}|_R}{\partial \dot q^k}}_{=: \heartsuit} \right) + \left( v_{Q_i}|_R \cdot \underbrace{\frac{\di}{\di t}\bigg|_R \frac{\partial v_{Q_i}|_R}{\partial \dot q^k} - \frac{\partial v_{Q_i}|_R}{\partial q^k}}_{=:\diamondsuit} \right) \right)
    \end{align*}
    Nel terzo passaggio abbiamo valutato le funzioni su $\hat\Gamma$ e usato la regola di Leibniz per la derivata del prodotto scalare.\\
    Per calcolare $\heartsuit$ usiamo la formula della velocità dell'osservazione di sopra 
    \[ \frac{\partial v_{Q_i}|_R}{\partial \dot q^k} = \pi_i\left( \frac{\partial \psi|_{\Sigma_t^N}^{-1}}{\partial q^k} \right) \]
    E allo stesso modo per $\diamondsuit$:
    \begin{align*}
        \frac{\di}{\di t}\bigg|_R \frac{\partial v_{Q_i}|_R}{\partial \dot q^k} - \frac{\partial v_{Q_i}|_R}{\partial q^k} & = \pi_i\left( \frac{\di}{\di t}\bigg|_R \frac{\partial \psi|_{\Sigma_t^N}^{-1}}{\partial q^k} - \frac{\partial}{\partial q}\frac{\partial }{\partial t}\bigg|_R \psi|_{\Sigma_t^N}^{-1} -\sum_l \frac{\partial^2 \psi|_{\Sigma_t^N}^{-1}}{\partial q^k \partial q^l} \frac{\di q^l}{\di t} \right) \\
        & = \pi_i\left( \frac{\di}{\di t}\bigg|_R \frac{\partial \psi|_{\Sigma_t^N}^{-1}}{\partial q^k} + \sum_l \frac{\partial^2 \psi|_{\Sigma_t^N}^{-1}}{\partial q^l \partial q^k} \frac{\di q^l}{\di t} - \frac{\partial}{\partial q}\frac{\partial }{\partial t}\bigg|_R \psi|_{\Sigma_t^N}^{-1} -\sum_l \frac{\partial^2 \psi|_{\Sigma_t^N}^{-1}}{\partial q^k \partial q^l} \frac{\di q^l}{\di t} = 0 \right)
    \end{align*}
    Dove i termini si annullano per il teorema di Schwarz e il fatto che $R$ non dipende da $q^k$.
    \qed
\end{proposition}

Nel caso in cui $\Gamma: I\to\V^{n+1}$ non si fattorizzi attraverso una singola carta locale, si può scegliere un ricoprimento aperto di $I$ con sottointervalli abbastanza piccoli e lavorare localmente su quello.

\newpage
\section{Equazioni di Eulero--Lagrange}

Adesso deriveremo e studieremo le equazioni cardinali della meccanica Lagrangiana, le equazioni di Eulero--Lagrange, che governano un sistema vincolato sotto le ipotesi che i vincoli siano olonomi e ideali.

\begin{theorem}{Equazioni di Eulero--Lagrange}{Eulero--Lagrange 1}
    Sia $\{Q_1,\ ...\ ,\ Q_N\}$ un sistema di $N\ge 1$ punti materiali sottoposto a un sistema $\{f_1,\ ...\ ,\ f_C\}$ di $0\le C\le 3N$ vincoli olonomi. Sia inoltre $\Gamma:I\to\V^{n+1}\subset\V^{3N+1}$ una linea di universo nello spaziotempo delle configurazioni del sistema vincolato con $I\subset E^1$ abbastanza piccolo da ammettere una rappresentazione in coordinate locali naturali $\Gamma(t) = \psi^{-1}(t,\bar q(t))$.\\
    Se $\Gamma$ soddisfa le equazioni di Newton 
    \[ m_{Q_i}a_{Q_i}|_R(t) = F_{R,i}(\Gamma(t),v_R(t)) + \Phi_i \]
    Dove $\Phi_i$ sono reazioni vincolari che verificano la proprietà di vincoli ideali, allora soddisfa anche le \bemph{equazioni di Eulero--Lagrange}:
    \[ \frac{\di}{\di t}\left[ \frac{\partial \T_R}{\partial \dot q^k} \right]_{\hat\Gamma(t)} - \left[ \frac{\partial\T_R}{\partial q^k} \right]_{\hat\Gamma(t)} = [\mc Q_{R,k}]_{\hat\Gamma(t)} \]
    Dove 
    \[ \mc Q_{R,k}(t,\bar q, \dot{\bar q}) = \left( F_R \cdot \frac{ \partial \psi|_{\Sigma_t^N}}{\partial q^k} \right) \]
    sono le componenti della forza $F_R = (F_{R,1},\ ...\ ,\ F_{R,N}) \in V_t^N$ lungo i vettori tangenti di $T_q\Q_t$.
    \proof 
    Scriviamo il sistema delle equazioni di Newton come un'unica equazione 
    \[ \begin{pmatrix}
        m_{Q_1}a_{Q_1}|_R \\ \vdots \\ m_{Q_N}a_{Q_N}|_R
    \end{pmatrix} = \begin{pmatrix}
        F_{R,1} \\ \vdots \\ F_{R,N}
    \end{pmatrix} + \begin{pmatrix}
        \Phi_1 \\ \vdots \\ \Phi_N
    \end{pmatrix} = F_R + \Phi \]
    Facendo il prodotto scalare con i vettori tangenti da entrambi i lati concludiamo immediatamente, ricordando le definizioni di vincoli ideali e l'ultima proposizione della sezione precedente
    \[ \underbrace{ \sum_i m_{Q_i} \left( a_{Q_i}|_R(t)\cdot \pi_i\left( \frac{\partial \psi|_{\Sigma_t^N}^{-1}(t,\bar q,\dot{\bar q})}{\partial q^k} \right) \right) }_{ = \frac{\di}{\di t}\left[ \frac{\partial \T_R}{\partial \dot q^k} \right]_{\hat\Gamma(t)} - \left[ \frac{\partial\T_R}{\partial q^k} \right]_{\hat\Gamma(t)} } = \underbrace{ \left( F_R \cdot \frac{ \partial \psi|_{\Sigma_t^N}}{\partial q^k} \right) }_{= [\mc Q_{R,k}]_{\hat\Gamma(t)}} + \underbrace{ \left( \Phi \cdot \frac{ \partial \psi|_{\Sigma_t^N}}{\partial q^k} \right)}_{= 0}  \]
    \qed 
\end{theorem}

Notiamo che il vantaggio delle equazioni di Eulero--Lagrange è che grazie all'ipotesi di idealità dei vincoli, queste non dipendono dalle reazioni vincolari, che nelle equazioni di Newton appaiono come incognite.\\
Tuttavia il teorema di cui sopra non risulta essere davvero soddisfacente: non ci dice che possiamo trovare una soluzione\footnote{Che in linea di massa potrebbe anche non esistere per quanto ne sappiamo ora} delle equazioni di Eulero--Lagrange ed essere certi che questa sia effettivamente una linea di universo che soddisfa le equazioni di Newton. Per ottenere questo risultato dobbiamo lavorare ancora un po'.\\
Ricordiamo che l'energia cinetica in coordinate naturali ha forma 
\[ \T_R = \sum_{k,l}A_{k,l}\dot q^k \dot q^l + \sum_k B_k \dot q^k + C  \]
Ciò implica che le equazioni di Eulero--Lagrange possono essere riscritte come 
\[ \sum_l A_{k,l}(t,\bar q(t))\frac{\di^2 q^l(t)}{\di t^2} = G_k\left(t, \bar q(t), \frac{\di\bar q(t)}{\di t}\right) \]
Per applicare il teorema \cref{th:esistenza e unicità locale 1} dobbiamo dimostrare che la matrice $A(t,\bar q) = (A_{k,l}(t,\bar q))_{k,l}$ è invertibile, in modo da poter riscrivere l'equazione in forma normale.

\begin{theorem}{Invertibilità della matrice per l'energia cinetica}{invertibilità della matrice per l'energia cinetica}
    Nel contesto del teorema precedente, la matrice $A(t,\bar q)$ è definita positiva per ogni $(t,\bar q)$ e dunque è invertibile.
    \proof 
    Ricordiamo che una matrice è definita positiva se e solo se $\bar v^T A \bar v > 0$ per ogni vettore $\bar v$  non nullo; in particolare possiamo riscrivere questa condizione per esteso usando la formula per $A_{k,l}$ come 
    \begin{align*} \sum_{k,l} A_{k,l} v^k v^l & = \sum_i \frac{m_{Q_i}}{2} \left( \pi_i\left( \sum_k v^k \frac{\partial \psi|_{\Sigma_t^N}^{-1}}{\partial q^k} \right) \cdot \pi_i\left( \sum_k v^k \frac{\partial \psi|_{\Sigma_t^N}^{-1}}{\partial q^k} \right) \right) \\ & = \sum_i \frac{m_{Q_i}}{2} \left| \pi_i\left( \sum_k v^k \frac{\partial \psi|_{\Sigma_t^N}^{-1}}{\partial q^k} \right) \right|^2 > 0 \end{align*}
    Questa espressione è sicuramente maggiore o uguale a zero, è una somma di termini non negativi (le norme) scalati di costanti positive (le masse); potrebbe essere nulla, ma grazie al fatto che le derivate di $\psi|_{\Sigma_t^N}^{-1}$ sono una base di $T_q\Q_t$, un vettore ha componente nulla rispetto a ogni elemento della base se e solo se è il vettore nullo, dunque concludiamo osservando che $A$ essendo simmetrica e definita positiva è invertibile.
    \qed
\end{theorem}

\begin{corollary}{Esistenza e unicità locale per le equazioni di Eulero--Lagrange}{esistenza e unicità locale 2}
    Le equazioni di Eulero--Lagrange possono essere riscritte in forma normale come
    \[ \frac{\di^2 q^l(t)}{\di t^2} =  \sum_k \frac{1}{2} (A^{k,l}(t,\bar q(t)))^{-1} G_k\left(t, \bar q(t), \frac{\di\bar q(t)}{\di t}\right) \] 
    E dunque per ogni dato iniziale $(t_0,\bar q_0, \dot{\bar q}_0)$ esiste ed è unica la soluzione locale.
\end{corollary}

Questo teorema ci dice qualcosa di molto forte, e in un certo senso "giustifica" il passaggio dalla meccanica Newtoniana a quella Lagrangiana: la seconda è in grado di darci delle equazioni del moto deterministiche sia nei sistemi svincolati, dove era sufficiente quella Newtoniana, sia in quelli vincolati sotto opportune ipotesi di idealità.\\
In molti casi pratici, ma non in tutti, le forze $F_R = (F_{R,1},\ ...\ ,\  F_{R,N})$ hanno una forma particolare, ovvero ammettono un cosiddetto \bemph{potenziale}: ciò significa che esiste una funzione liscia $\U : \V^{3N+1}\to\R$ tale che per ogni punto $p \in \V^{3N+1}$ e ogni vettore $v \in V_{T(p)}^N$ valga
\[ F_{R}(p) \cdot v = -\frac{\di}{\di\varepsilon}\U(p+\varepsilon v)\bigg|_{\varepsilon = 0} \]
Che in coordinate standard adattate $(t,\bar{\bar x})$ assume la forma familiare
\[ \bar F_{R,i}(t,\bar{\bar x}) = - \nabla_{\bar x_i} \U(t, \bar{\bar x}) \]
Dove il segno negativo è per convenzione. In questo caso, le componenti Lagrangiane delle forze sono date da
\[ Q_{R,k}(t,\bar q) = F_R \cdot \frac{\partial \psi|_{\Sigma_t^N}^{-1}(\bar q)}{\partial q^k} = -\frac{\partial\U(t,\bar q)}{\partial q^k}\]
Dove con un piccolo abuso di notazione abbiamo indicato $\U(t,\bar q) = \U(\psi^{-1}(t,\bar q))$ restringendoci a una carta locale su $\V^{n+1}$.

\begin{corollary}{Equazioni di Eulero--Lagrange con un potenziale}{Eulero--Lagrange 2}
    Nel contesto dei teoremi precedenti, si supponga che le forze $F_R$ ammettano un potenziale $\U$; definiamo allora in coordinate locali naturali la \bemph{Lagrangiana} del sistema $\L_R : A(\V^{n+1})\to \R$ come
    \[ \L_R(t,\bar q, \dot{\bar q}) := \T_R(t,\bar q, \dot{\bar q}) - \U(t, \bar q) \]
    Sotto questa ipotesi, le equazioni di Eulero--Lagrange possono essere riscritte come
    \[ \frac{\di}{\di t}\left[ \frac{\partial \L_R}{\partial \dot q^k} \right]_{\hat\Gamma(t)} - \left[ \frac{\partial \L_R}{\partial q^k} \right]_{\hat\Gamma(t)} = 0\]
    \proof 
    Notando che $\U$ non dipende dalle coordinate puntate, il primo termine si riconosce essere semplicemente la derivata dell'energia cinetica $\T_R$, mentre il secondo termine per linearità è
    \[ \left[ \frac{\partial \L_R}{\partial q^k} \right]_{\hat\Gamma(t)} = \left[ \frac{\partial \T_R}{\partial q^k} \right]_{\hat\Gamma(t)} - \underbrace{\left[ \frac{\partial \U_R}{\partial q^k} \right]_{\hat\Gamma(t)}}_{= - [Q_{R,k}]_{\hat\Gamma(t)}} \]
    Poi basta riordinare per recuperare le equazioni di Eulero--Lagrange.
\end{corollary}

Nel contesto della meccanica Lagrangiana, spesso un sistema è descritto puramente in termini della sua Lagrangiana $\L_R : A(\V^{n+1})\to\R$ senza specificare le forze, anche se questa non dovesse assumere la forma di una Lagrangiana "standard" $\T_R -\U$, ad esempio lasciando che un potenziale dipenda anche dalle coordinate puntate.\\
Notiamo che due Lagrangiane diverse possono generare le stesse equazioni di Eulero--Lagrange: per una mappa liscia $g:\V^{n+1}\to\R$ definiamo la sua derivata totale formale come una funzione $\di_t g: A(\V^{n+1})\to\R$ in coordinate naturali come
\[ \di_t g(t,\bar q, \dot{\bar q}) := \frac{\partial g(t,\bar q)}{\partial t} + \sum_{k} \frac{\partial g(t,\bar q)}{\partial q^k} \]
Allora per calcolo diretto si verifica che $\L_R$ e $\L_R +\di_t g$ producono le stesse equazioni di Eulero--Lagrange; per questo motivo a volte è utile pensare alle Lagrangiane come a classi di equivalenza $[\L]$ rispetto alla relazione $\L \sim \L + \di_t g$.

\chapter{Simmetrie e leggi di conservazione}

In questa parte studieremo degli interessanti legami tra le simmetrie di un sistema fisico descritto da una Lagrangiana $\L: A(\V^{n+1})\to\R$ e le quantità che rimangono costanti durante l'evoluzione dinamica del sistema. Partiremo da un caso molto semplice, quello delle coordinate cicliche, per poi studiare i teoremi di Noether e di Jacobi.

\newcommand\I{\mc I}
\section{Coordinate cicliche}

\begin{definition}{Coordinate cicliche}{coordinate cicliche}
    Si consideri un sistema fisico descritto da uno spaziotempo delle configurazioni $T:\V^{n+1}\to E^1$ e da una Lagrangiana $\L:A(\V^{n+1})\to\R$ e sia $(U,\psi)$ una carta locale per $\V^{n+1}$ con le corrispondenti coordinate locali naturali $(t, q, \dot q) := (t,\bar q, \dot{\bar q})$.\\
    Una coordinata $q^j$ si dice \bemph{ciclica} se la derivata parziale 
    \[ \frac{\partial \L(t,q,\dot q)}{\partial q^j} = 0 \]
    È nulla per tutti i punti di $\psi(U)\times\R^n$.
\end{definition}

Notiamo che il significato geometrico di questa richiesta è che la Lagrangiana sia invariante per traslazione lungo questa coordinata, ovvero che per ogni $b\in\R$ sufficientemente piccolo valga
\[\L(t,q^1,\ ...\ ,\ q^j,\ ...\ ,\ q^n,\dot q) = \L(t,q^1,\ ...\ ,\ q^j+b,\ ...\ ,\ q^n,\dot q)\]
Ne segue che le coordinate cicliche esprimono una simmetria (locale) per traslazione del sistema fisico.

\begin{proposition}{Momento coniugato}{momento coniugato}
    Nel contesto della definizione precedente, definendo il \bemph{momento coniugato associato} a una coordinata ciclica $q^j$ come 
    \[ p_j(t,q,\dot q) = \frac{\partial \L(t,q,\dot q)}{\partial \dot q^j} \]
    Si ha che questo è un \bemph{integrale primo} del moto, ovvero la sua valutazione lungo una linea di universo $\hat\Gamma(t)$ che risolve le equazioni di Eulero--Lagrange è costante rispetto al tempo.
    \proof 
    Richiamiamo le equazioni di Eulero--Lagrange per la coordinata ciclica $q^j$:
    \[ \frac{\di}{\di t} \underbrace{\left[ \frac{\partial \L}{\partial \dot q^j} \right]_{\hat\Gamma(t)}}_{=: [p_j]_{\hat\Gamma(t)}} - \underbrace{\left[ \frac{\partial \L_R}{\partial q^j} \right]_{\hat\Gamma(t)}}_{= 0} = 0\] 
    \qed
\end{proposition}

In linea di principio, l'integrale primo $p_j(t,q,\dot q)$ è definito solo localmente sullo spaziotempo degli atti di moto perchè dipende da coordinate locali; nel contesto del teorema di Noether, che occuperà tutta la prossima sezione, vedremo quando è possibile "incollare" integrali primi locali per ottenerne uno globale $\I : A(\V^{n+1})\to\R$.\\
È importante sottolineare che l'integrale primo è costante nel tempo \textit{per una linea di universo fissata}, in generale linee di universo con dati iniziali diversi possono avere integrali primi anche molto diversi.\\
Ovviamente, più coordinate cicliche generano più integrali primi del moto.

\begin{example}{Punti materiali svincolati}{punti materiali svincolati}
    Consideriamo un sistema di $N\ge 1$ punti materiali svincolati e senza forze in un riferimento inerziale $R$ con coordinate standard $(t,\bar{\bar x}) \in \R^{3N+1}$.\\
    La Lagrangiana di questo sistema è data da
    \[ \L(t,\bar{\bar x},\dot{\bar{\bar x}})  = \sum_i \frac{m_{Q_i}}{2}\| \dot{\bar x}_i \|^2\]
    Osserviamo che tutte le coordinate $x_i^k$ con $i \le N$ e $k\le 3$ sono cicliche; i momenti coniugati sono dati da
    \[ p_k^i(t,\bar{\bar x}, \dot{\bar{\bar x}}) = \frac{\partial \L(t,\bar{\bar x},\dot{\bar{\bar x}})}{\partial \dot x_k^i} = m_{Q_i}\dot x_i^k \]
    E, per la proposizione precedente, definiscono degli integrali primi del moto.\\
    Si nota immediatamente che questi non sono nient'altro che le componenti degli impulsi del sistema, quindi ogni impulso $\bar p^i$ è conservato durante il moto del sistema
\end{example}

\begin{example}{Punti materiali svincolati con forze dipendenti dalle posizioni relative}{Punti materiali svincolati con forze dipendenti dalle posizioni relative}
    Supponiamo di introdurre nel sistema precedente delle forze che derivano da un potenziale $\U$ dipendente soltanto dalle differenze $\bar x_i -\bar x_j$ tra le posizioni: la Lagrangiana allora diventa 
    \[ \L(t,\bar{\bar x},\dot{\bar{\bar x}})  = \sum_i \frac{m_{Q_i}}{2}\| \dot{\bar x}_i \|^2 - \U(\bar x_1 - \bar x_2,\ ...\ ,\  \bar x_{N-1} - \bar x_N)\]
    Dove abbiamo scritto $\U$ puramente in funzione delle differenze tra indici consecutivi\footnote{Perchè possiamo farlo senza perdita di generalità? Lo lasciamo come esercizio per il lettore}.\\
    Almeno così ad occhio, non sembrano esserci coordinate cicliche... o invece ce ne sono? Facciamo un cambiamento di coordinate: 
    \[ \bar X := \frac{ \sum_i m_{Q_i}\bar x_i }{\sum_i m_{Q_i}} \]
    \[ \bar\xi_i := \bar x_i - \frac{ \sum_{j = i+1} m_{Q_j}\bar x_j }{\sum_{j=i+1} m_{Q_j}} \]
    Dette rispettivamente \bemph{centro di massa} e \bemph{coordinate di Jacobi} del sistema.\\
    Si può dimostrare per induzione su $N$ che la Lagrangiana di questo sistema è scritta come 
    \[ \L\left(t,\bar X, \bar{\bar \xi}, \dot{\bar X}, \dot{\bar{\bar\xi}}\ \right) = \frac{M}{2}\| \dot{\bar X} \|^2 + \sum_i \frac{\mu_i}{2} \|\dot{\bar\xi}_i\|^2 - \tilde\U(\bar{\bar\xi}) \]
    Dove $M$ è la massa totale e le $\mu_i$ sono dette \bemph{masse ridotte} e sono costanti definite in termini delle masse un modo che non ci interessa.\\
    Poichè il potenziale dipende solo dalle $\bar\xi_i$, le coordinate di $\bar X=(X^1,X^2,X^3)$ sono coordinate cicliche e gli integrali primi di questo moto sono quelle che riconosciamo essere le componenti dell'impulso totale 
    \[ [\bar p]_{\hat\Gamma(t)} = \sum_i m_{Q_i} \frac{\di \bar x_i(t)}{\di t} \]
\end{example}

Questo esempio ci illustra come a volte le coordinate cicliche siano nascoste e possano essere recuperate solo attraverso cambiamenti di coordinate anche piuttosto elaborati.\\
Questo diventerà molto più chiaro nel contesto generale del teorema di Noether, che illustra esattamente le corrispondenze tra le simmetrie del sistema fisico e le quantità conservate.

\newpage
\newcommand\loc{\text{loc}}
\section{Il Teorema di Noether}

Per formulare e dimostrare il teorema di Noether dobbiamo considerare una classe particolare di automorfismi dello spaziotempo delle configurazioni, ovvero quelli che non trasformano il tempo.

\begin{definition}{Automorfismo che preserva le fibre}{automorfismo che preserva le fibre}
    Un \bemph{automorfismo che preserva le fibre} dello spaziotempo delle configurazioni $T:\V^{n+1}\to E^1$ è un automorfismo $F:\V^{n+1}\iso \V^{n+1}$ tale che $T = T \circ F = T \circ F^{-1}$, ovvero tale che il seguente diagramma commuti:
    \[\begin{tikzcd}
    	{\V^{n+1}} && {\V^{n+1}} \\
	    \\
	    & {E^1}
	    \arrow["F"{description}, from=1-1, to=1-3]
	    \arrow["T"{description}, from=1-1, to=3-2]
	    \arrow["T"{description}, from=1-3, to=3-2]
    \end{tikzcd}\]
    Un \bemph{gruppo a un parametro di automorfismi che preservano le fibre} è una mappa liscia $F:\R\times\V^{n+1}\to\V^{n+1}$ tale che la mappa $F_\bullet : s \mapsto F_s : \V^{n+1}\iso\V^{n+1}$ sia un omomorfismo di gruppi dal gruppo additivo dei numeri reali al sottogruppo che preserva le fibre del gruppo degli automorfismi di $\V^{n+1}$, ovvero per ogni $s,r \in \R$ valgano:\begin{enumerate}
        \item $T = T \circ F_s$
        \item $F_0 = \id_{\V^{n+1}} $
        \item $F_{s+r} = F_s\circ F_r$
    \end{enumerate}
\end{definition}

Curiosamente, l'immagine di $F_\bullet$ deve essere dunque un sottogruppo abeliano del gruppo degli automorfismi\footnote{E dunque, del sottogruppo di quelli che preservano le fibre} di $\V^{n+1}$.\\
Per facilitare i calcoli sarebbe bello riuscire a esprimere l'azione di questi automorfismi in coordinate, ma c'è un problema: per qualsiasi carta locale adattata $(U,\psi)$, non è assolutamente detto che $F_s(U) \subset U$, quindi che si fa?\\
È presto detto, basta restringersi a un intervallo $(-\varepsilon, \varepsilon)$ e un aperto $V\subset U$ sufficientemente piccoli da avere $F((-\varepsilon, \varepsilon)\times V)\subset U$. In questa situazione possiamo esprimere l'automorfismo in termini di coordinate locali naturali nel seguente diagramma:
\[\begin{tikzcd}
	& {(-\varepsilon,\varepsilon)\times V} && U & \\
	\\
	{(-\varepsilon,\varepsilon)\times \psi(V)} &&&& {\psi(U)} \\
	{(s,(t,q))} &&&& {(t,q'(s,t,q))}
	\arrow["F"{description}, from=1-2, to=1-4]
	\arrow["{\id_{(-\varepsilon, \varepsilon)}\times \psi}"{description}, from=1-2, to=3-1]
	\arrow["\psi"{description}, from=1-4, to=3-5]
	\arrow["{\psi\circ F\circ (\id_{(-\varepsilon, \varepsilon)}\times \psi)^{-1}}"{description}, from=3-1, to=3-5]
	\arrow[maps to, from=4-1, to=4-5]
\end{tikzcd}\]
Dove ovviamente $q = (q^1,\ ...\ ,\ q^n)$ e $q' = (q'^1,\ ...\ ,\ q'^n)$.\\
Possiamo estendere l'azione di $F$ allo spaziotempo degli atti di moto $A(\V^{n+1})$ mandando $(s,(t,q,\dot q))$ in $(t,q'(s,t,q),\dot q'(s,t,q,\dot q))$, dove 
\[ \dot q'^k (s,t,q,\dot q) = \frac{\partial q'(s,t,q)}{\partial t} + \sum_l \frac{\partial q'^k(s,t,q)}{\partial q^l} \dot q^l \]

\begin{definition}{Sistema invariante}{sistema invariante}
    Si consideri un sistema fisico con spaziotempo delle configurazioni $T:\V^{n+1}\to E^1$ descritto dalla Lagrangiana $\L:A(\V^{n+1})\to\R$.\\
    Questo sistema si dice \bemph{debolmente invariante} rispetto a un gruppo di automorfismi a un parametro $F:\R\times\V^{n+1}\to\V^{n+1}$ che preservano le fibre se, in ogni carta locale adattata $(U,\psi)$ e aperto $V\subset U$ sufficientemente piccolo come visto sopra, esiste una funzione liscia $g: \V^{n+1}\to\R$ tale che
    \[ \frac{\partial \L(t, q'(s,t,q), \dot q'(t,s,q,\dot q))}{\partial s}\bigg|_{s=0} = \di_t g(t,q,\dot q) := \frac{\partial g(t,q)}{\partial t} + \sum_k \frac{\partial g(t,q)}{\partial q^k}\dot q^k \]
    In particolare, se $g$ è la funzione costantemente nulla, il sistema si dice \bemph{invariante}.\\
    $F$ è detto \bemph{gruppo di simmetria (debole)} del sistema fisico.
\end{definition}

Chiaramente l'idea alla base dell'invarianza debole è l'osservazione fatta qualche pagina fa secondo cui le Lagrangiane $\L$ e $\L+\di_t g$ producono le stesse equazioni di Eulero--Lagrange: infatti, l'invarianza debole afferma che non la Lagrangiana $\L$ in sè ma la sua classe di equivalenza $[\L]$ è invariante rispetto all'azione di $F$.\\
Si può introdurre anche una nozione di invarianza (debole) locale quando si dispone non di un automorfismo globale da restringere opportunamente ma di una famiglia di diffeomorfismi locali $F_\loc : (-\varepsilon, \varepsilon)\times \psi(V)\to\psi(U)$; dato che comunque noi lavoreremo in carte locali, le dimostrazioni che seguono si applicano anche al caso in cui si disponga soltanto di una simmetria locale.

\begin{theorem}{Teorema di Noether}{teorema di Noether}
    Nel contesto della definizione precedente, le funzioni lisce $I_{F,g}:\psi(V)\times\R^n\to\R$ definite da
    \[ I_{F,g}(t,q,\dot q) = -g(t,q) + \sum_k \frac{\partial \L(t,q,\dot q)}{\partial \dot q^k}\frac{q'^k(s,t,q)}{\partial s}\bigg|_{s=0}  \]
    Sono integrali primi del moto, vale a dire che la loro valutazione lungo una linea di universo $\hat\Gamma:I\to A(\V^{n+1})$ che soddisfa le equazioni di Eulero--Lagrange è costante.\\
    Questi integrali primi si assemblano per incollamento in un integrale primo globale $\I_{F,g}:A(\V^{n+1})\to\R$ detto \bemph{integrale primo di Noether}.
    \proof 
    Analizziamo la condizione di invarianza debole 
    \[ \frac{\partial \L(t, q'(s,t,q), \dot q'(t,s,q,\dot q))}{\partial s}\bigg|_{s=0} = \di_t g(t,q,\dot q) \]
    Sviluppando il lato sinistro 
    \begin{align*}
        \frac{\partial \L(t, q'(s,t,q), \dot q'(t,s,q,\dot q))}{\partial s}\bigg|_{s=0} & = \sum_k \frac{\partial \L(t,q',\dot q' )}{\partial q'^k} \frac{\partial q'^k}{\partial s} \bigg|_{s=0} + \sum_k \frac{\partial \L(t,q',\dot q' )}{\partial \dot q'^k} \frac{\partial \dot q'^k}{\partial s} \bigg|_{s=0} \\ & = \sum_k \frac{\partial \L(t,q,\dot q )}{\partial q^k} \frac{\partial q'^k}{\partial s} \bigg|_{s=0} + \sum_k \frac{\partial \L(t,q,\dot q )}{\partial \dot q^k} \underbrace{\frac{\partial \dot q'^k}{\partial s} \bigg|_{s=0}}_{=: \heartsuit}
    \end{align*}
    Dove per passare alla seconda riga abbiamo usato il fatto che $q'|_{s=0} = q$ e $\dot q'|_{s=0} = \dot q$. A questo punto esaminando $\heartsuit$ vediamo che 
    \begin{align*}
        \frac{\partial \dot q'^k}{\partial s} \bigg|_{s=0} & = \frac{\partial }{\partial s}\left( \frac{\partial q'^k}{\partial t} + \sum_l \frac{\partial q'^k}{\partial q^l}\dot q^l \right)\bigg|_{s=0} \\ & = \left( \frac{\partial^2 q'^k }{\underbrace{\partial s \partial t}_{= \partial t \partial s}} + \sum_l \frac{\partial^2 q'^k}{\underbrace{\partial s \partial q^l}_{= \partial q^l \partial s}} \dot q^l \right)\bigg|_{s=0} \\ & = \frac{\partial }{\partial t}\left( \frac{\partial q'^k}{\partial s} \bigg|_{s=0} \right) + \sum_k \frac{\partial}{\partial q^l} \left( \frac{\partial q'^k}{\partial s}\bigg|_{s=0} \right)\dot q^l = \frac{\di}{\di t}\left( \frac{\partial q'^k}{\partial s}\bigg|_{s=0} \right)
    \end{align*}
    Adesso, ricordandoci che $\Gamma$ soddisfa le equazioni di Eulero--Lagrange 
    \[ \frac{\di}{\di t}\left[ \frac{\partial \L}{\partial \dot q^k} \right]_{\hat\Gamma(t)} =  \left[ \frac{\partial \L}{\partial q^k} \right]_{\hat\Gamma(t)} \]
    Possiamo scrivere 
    \begin{align*}
        \left[ \frac{\partial \L(t,q',\dot q')}{\partial s}\bigg|_{s=0} \right]_{\hat\Gamma(t)} & = \sum_k \left( \left[ \frac{\partial \L}{\partial q^k} \right]_{\hat\Gamma(t)}\left[ \frac{\partial q'^k}{\partial s}\bigg|_{s=0} \right]_{\hat\Gamma(t)} + \left[ \frac{\partial \L}{\partial \dot q^k} \right]_{\hat\Gamma(t)}\frac{\di}{\di t} \left[ \frac{\partial q'^k}{\partial s}\bigg|_{s=0} \right]_{\hat\Gamma(t)} \right) \\
        & = \sum_k \left( \frac{\di}{\di t}\left[ \frac{\partial \L}{\partial \dot q^k} \right]_{\hat\Gamma(t)}\left[ \frac{\partial q'^k}{\partial s}\bigg|_{s=0} \right]_{\hat\Gamma(t)} + \left[ \frac{\partial \L}{\partial \dot q^k} \right]_{\hat\Gamma(t)}\frac{\di}{\di t} \left[ \frac{\partial q'^k}{\partial s}\bigg|_{s=0} \right]_{\hat\Gamma(t)} \right) \\ 
        & = \frac{\di}{\di t}\left[ \sum_k \frac{\partial \L}{\partial\dot q^k} \frac{\partial q'^k}{\partial s}\bigg|_{s=0} \right]_{\hat\Gamma(t)} = \frac{\di}{\di t}[g]_{\hat\Gamma(t)}
    \end{align*}
    Portando tutto al lato sinistro ne segue che 
    \[ \frac{\di}{\di t} [I_{F,g}]_{\hat\Gamma(t)} = \frac{\di}{\di t}\left[ -g + \sum_k \frac{\partial \L}{\partial \dot q^k}\frac{\partial q'^k}{\partial s}\bigg|_{s=0} \right]_{\hat\Gamma(t)} = 0 \]
    Che implica, dato che $I$ è connesso, che $[I_{F,g}]_{\hat\Gamma(t)}$ sia costante nel tempo.\\
    Per dimostrare che possiamo incollare gli integrali locali dobbiamo mostrare che, date due carte $(U,\psi)$ e $(\tilde U, \tilde \psi)$, su $\varnothing \neq V\cap \tilde V \subset U\cap \tilde U$ valga 
    \[ I_{F,g}(t,q,\dot q) = \left[ -g(t,q) + \sum_k \frac{\partial \L(t,q,\dot q)}{\partial \dot q^k}\frac{\partial q'^k}{\partial s}\bigg|_{s=0} \right]_{\hat\Gamma(t)} = \left[ -g(t,\tilde q) + \sum_k \frac{\partial \L(t,\tilde q,\dot{\tilde q})}{\partial \dot{\tilde q}^k}\frac{\partial \tilde q'^k}{\partial s}\bigg|_{s=0} \right]_{\hat\Gamma(t)} = I_{F,g}(t,\tilde q,\dot {\tilde q}) \]
    Per i termini che coinvolgono $g$ questo è automatico perchè richiediamo che sia definita globalmente; vediamo ora che 
    \begin{align*}
        \sum_k \frac{\partial \L(t,q,\dot q)}{\partial \dot q^k}\frac{\partial q'^k}{\partial s}\bigg|_{s=0} & = \sum_{k,l} \frac{\partial \L(t,q,\dot q)}{\partial \dot{\tilde q}^l} \frac{\partial \dot{\tilde q}^l}{\partial \dot q^k} \frac{\partial q'^k}{\partial s}\bigg|_{s=0} \\ & = \sum_{k,l} \frac{\partial \L(t,q,\dot q)}{\partial \dot{\tilde q}^l} \frac{\partial \tilde q^l}{\partial q^k} \frac{\partial q'^k}{\partial s}\bigg|_{s=0} = \sum_l \frac{\partial \L}{\partial \dot{\tilde q }^l }\frac{ \partial \tilde q'^l}{\partial s}\bigg|_{s=0}
    \end{align*}
    Dove l'ultimo passaggio si legge più in dettaglio come 
    \[ \frac{\partial \tilde q'^l}{\partial s}\bigg|_{s=0} = \frac{\partial \tilde q^l(t, q'(s,t,q))}{\partial s}\bigg|_{s=0} = \sum_k \left( \frac{\partial \tilde q^l}{\partial q'^k}\frac{\partial q'^k}{\partial s}\bigg|_{s=0} \right) = \sum_k \frac{\partial \tilde q^l}{\partial q^k} \left( \frac{\partial q'^k}{\partial s}\bigg|_{s=0} \right)\]
    \qed
\end{theorem}

\newpage
\renewcommand\H{\mc H}
\section{Il teorema di Jacobi}

In questa parte introduciamo il teorema di Jacobi, che è complementare al teorema di Noether perchè si occupa dell'invarianza rispetto alle traslazioni temporali invece che alle variazioni lungo una coordinata spaziale.

\begin{theorem}{Teorema di Jacobi}{teorema di Jacobi}
    Si consideri un sistema fisico $T:\V^{n+1}\to E^1$ descritto da una Lagrangiana $\L : A(\V^{n+1})\to\R$.\\
    Data una carta locale adattata $(U,\psi)$, definiamo la \bemph{funzione Hamiltoniana locale} in $(U,\psi)$ del sistema come la funzione liscia $\H : \psi(U)\times\R^n\to \R$ data da 
    \[ \H(t,q,\dot q) = \sum_k \dot q^k\frac{\partial \L(t,q,\dot q)}{\partial \dot q^k} - \L(t,q,\dot q) \]
    Allora lungo ogni linea di universo $\Gamma:I\to\V^{n+1}$ che risolve le equazioni di Eulero--Lagrange vale l'identità 
    \[ \frac{\di}{\di t} [\H]_{\hat\Gamma(t)} = -\left[ \frac{\partial \L}{\partial t} \right]_{\hat\Gamma(t)} \]
    \proof 
    La dimostrazione procede per calcolo diretto 
    \begin{align*}
        \frac{\di}{\di t}[\H]_{\hat\Gamma(t)} & = \frac{\di}{\di t}\left[ \sum_k \dot q^k \frac{\partial \L}{\partial \dot q^k} - \L \right]_{\hat\Gamma(t)} \\
        & = \frac{\di}{\di t}\left( \frac{\di q^k}{\di t}\left[ \frac{\partial \L}{\partial \dot q^k} \right]_{\hat\Gamma(t)} - [\L]_{\hat\Gamma(t)}  \right) \\ 
        & = \sum_k \left( \frac{\di^2 q^k}{\di t^2}\left[ \frac{\partial \L}{\partial \dot q^k} \right]_{\hat\Gamma(t)} + \frac{\di q^k}{\di t} \underbrace{\frac{\di}{\di t}\left[ \frac{\partial \L}{\partial \dot q^k} \right]_{\hat\Gamma(t)}}_{\text{Eulero--Lagrange}} \right) - \frac{\di}{\di t}[\L]_{\hat\Gamma(t)} \\ 
        & = \sum_k\left( \underbrace{\frac{\di^2 q^k}{\di t^2}\left[ \frac{\partial \L}{\partial \dot q^k} \right]_{\hat\Gamma(t)}}_{=:\heartsuit_k } + \underbrace{\frac{\di q^k}{\di t} \left[ \frac{\partial \L}{\partial q^k} \right]_{\hat\Gamma(t)}}_{=:\diamondsuit_k} \right) - \underbrace{\frac{\di}{\di t}[\L]_{\hat\Gamma(t)}}_{=:\clubsuit}
    \end{align*}
    Studiando $\clubsuit$ tramite la regola della catena otteniamo 
    \[ \clubsuit = \frac{\di}{\di t}\L\left(t,q(t)\dot q(t)\right) = \frac{\partial \L}{\partial t} + \sum_k\left( \underbrace{\frac{\di^2 q^k}{\di t^2}\left[ \frac{\partial \L}{\partial \dot q^k} \right]_{\hat\Gamma(t)}}_{=\heartsuit_k } + \underbrace{\frac{\di q^k}{\di t} \left[ \frac{\partial \L}{\partial q^k} \right]_{\hat\Gamma(t)}}_{=\diamondsuit_k} \right) \]
    Vediamo dunque che i termini si cancellano
    \[ \frac{\di  }{\di t}[\H]_{\hat\Gamma(t)} = \sum_k (\heartsuit_k + \diamondsuit_k) -\left[ \frac{\partial \L}{\partial t} \right]_{\hat\Gamma(t)} -\sum_k(\heartsuit_k + \diamondsuit_k) = -\left[ \frac{\partial \L}{\partial t} \right]_{\hat\Gamma(t)} \]
    \qed
\end{theorem}

\begin{corollary}{}{}
    Nella situazione del teorema precedente, se la Lagrangiana è costante nel tempo in coordinate locali naturali, allora l'Hamiltoniana è un integrale primo del moto
\end{corollary}

È chiaro che questo corollario dunque si riferisce a una simmetria del sistema, come il teorema di Noether di Noether, dato che possiamo interpretare la costanza nel tempo come un'invarianza del sistema rispetto a traslazioni temporali $(t,q,\dot q)\mapsto (t+\tau, q, \dot q)$.\\
È importante sottolineare che l'Hamiltoniana è una funzione \textit{locale}, che \textit{non} può essere estesa a una funzione globale, così come il suo integrale primo: infatti ricordando che una trasformazione di coordinate $(t,q,\dot q)\mapsto (t', q', \dot q')$ è data da 
\[ (t', q', \dot q') = \left( t', q'(t,q), \frac{\partial q'(t,q)}{\partial t} + \sum_k \frac{\partial q'^k(t,q)}{\partial q^k}\dot q'^k \right) \]
L'Hamiltoniana si trasforma 
\begin{align*}
    \H(t,q,\dot q) & = \sum_k \dot q^k\frac{\partial \L(t,q,\dot q)}{\partial \dot q^k} - \L(t,q,\dot q) \\
    & = \sum_k \dot q^k\frac{\partial \L(t',q',\dot q')}{\partial \dot q^k} - \L(t',q',\dot q')\\
    & = \sum_l \underbrace{\sum_k \dot q^k\frac{\partial \dot q'^l}{\partial \dot q^k} }_{= \dot q'^l - \frac{\partial q'^l}{\partial t}}\frac{\partial \L(t',q',\dot q')}{\partial \dot q'^k} - \L(t',q',\dot q') \\
    & = \H(t',q',\dot q') - \sum_l \frac{\partial q'^l}{\partial t}\frac{\partial \L(t',q',\dot q')}{\partial q'^l} \neq \H(t',q',\dot q')
\end{align*}
Dunque vediamo che, analogamente alle coordinate cicliche, l'integrale primo di Jacobi dipende dalla scelta di una carta.\\
Consideriamo un caso particolare, in cui la Lagrangiana ha forma standard $\L(t,q,\dot q) = \T_R(q,\dot q)-\U(q)$ ed è indipendente dal tempo $t$, con energia cinetica di forma 
\[ \T_R = \sum_{k,l} A_{k,l}(q)\dot q^k\dot q^l  \]
Allora l'Hamiltoniana locale è data dall'energia totale del sistema:
\begin{align*}
    \H(q,\dot q) & = \sum_k \dot q^k \frac{\partial \L(q,\dot q)}{\partial \dot q^k} - \L(q,\dot q)\\
    & = \sum_{k,l} 2 A_{k,l}(q)\dot q^k \dot q^l - \sum_{k,l} A_{k,l}(q)\dot q^k\dot q^l + \U(q) = \T_R(q,\dot q) + \U(q)
\end{align*}
In questa situazione, il teorema di Jacobi implica che localmente l'energia totale si conserva durante l'evoluzione dinamica del sistema.

\chapter{Meccanica Hamiltoniana}

In questa parte studieremo una formulazione alternativa della meccanica classica, ovvero la meccanica Hamiltoniana.\\
Sebbene sia equivalente alla formulazione Lagrangiana, presenta alcune caratteristiche interessanti:\begin{itemize}
    \item Le equazioni di Hamilton assumono una forma più simmetrica rispetto a quella delle equazioni di Eulero--Lagrange e possono essere scritte in un'utile forma algebrica attraverso la cosiddetta parentesi di Poisson che vedremo più avanti.
    \item Per chi sarà interessato a proseguire con un percorso legato alla fisica matematica, la meccanica Hamiltoniana è il punto di partenza per le formulazioni della meccanica statistica e la meccanica quantistica, che tuttavia sono al di fuori di questo corso.
\end{itemize}

\newcommand\Leg{\mathbb L}
\section{Spaziotempo delle fasi e trasformata di Legendre}

Mentre la meccanica Lagrangiana è stata formulata sullo spaziotempo degli atti di moto $T:A(\V^{n+1})\to E^1$ con coordinate locali naturali $(t,q,\dot q)$, nella meccanica Hamiltoniana si sostituiscono le coordinate puntate $\dot q = (\dot q^1,\ ...\ ,\  \dot q^n)$ coi cosiddetti momenti coniugati $p = (p_1,\ ...\ ,\ p_k)$.\\
Se si viene dalla meccanica Lagrangiana, questi sono definiti in termini della Lagrangiana $\L:A(\V^{n+1})\to\R$ tramite la formula 
\[ p_k = \frac{\partial \L(t,q,\dot q)}{\partial \dot q^k} \]
Ma è importante sottolineare che nella meccanica Hamiltoniana questi non sono definiti in termini delle coordinate puntate, ma sono vere coordinate indipendenti che le sostituiscono.\\
Ricordandoci la legge di trasformazione per le coordinate puntate $(t,q,\dot q)\mapsto(t', q', \dot q')$, ovvero 
\[ t' = t',\quad q'^k = q'^k(t,q),\quad \dot q'^k = \dot q'^k(t,q,\dot q) = \frac{\partial q'^k(t,q)}{\partial t} + \sum_l \frac{\partial q'^k(t,q) }{\partial q^l}\dot q^l  \]
Segue che i momenti coniugati devono trasformarsi nel seguente modo:
\[p'_k = \frac{\partial \L(t,q,\dot q)}{\partial \dot q^k} = \sum_l \frac{\partial \dot q^l}{\partial \dot q'^l} \frac{\partial \L(t,q,\dot q)}{\partial \dot q^l} = \sum_l \frac{\partial \dot q^l}{\partial \dot q'^l} p_l \]

\begin{definition}{Spaziotempo delle fasi}{spaziotempo delle fasi}
    Sia $T:\V^{n+1}\to E^1$ uno spaziotempo delle configurazioni per un sistema fisico con vincoli olonomi e sia $\A = \{ (U_i,\psi_i) \}_i$ l'atlante dato dalle carte locali adattate alla sommersione suriettiva $T$.\\
    Il corrispettivo \bemph{spaziotempo delle fasi} $T:F(\V^{n+1})\to E^1$ si definisce applicando la costruzione di fibrato jet alle carte di coordinate locali $\psi_i(U_i)\times\R^n$ e le mappe di trasferimento $\psi_j(U_i\cap U_j)\times \R^n \iso \psi_i(U_i\cap U_j)\times \R^n$ che cambiano le coordinate $(t,q,p)\mapsto (t',q',p')$ date da 
    \[ t'=t,\quad q'^k=q'^k(t,q),\quad p'_k = \sum_l \frac{\partial q^l(t,q)}{\partial q'^k}p_l  \]
    Con la sommersione suriettiva $(t,q,p)\mapsto t$
\end{definition}

Notiamo che richiedendo che i momenti coniugati rispondano alla stessa legge di trasformazione derivata dal formalismo Lagrangiano, otteniamo un'identità che coinvolge l'inversa della matrice Jacobiana del cambiamento di coordinate $q'=q'(t,q)$, definita dalle derivate di $q^l$ rispetto a $q'^k$.\\
Questo implica che i momenti coniugati $p_k$ si comportino \textit{dualmente} rispetto alle velocità $\dot q^k$, cosa che nel linguaggio della geometria differenziale viene espressa dicendo che i $p_k$ sono \bemph{vettori cotangenti} alle coordinate $q^k$\\
La formula che determina i momenti coniugati a partire dalla Lagrangiana definisce un'importante trasformazione globale che ci permette di passare dalla meccanica Lagrangiana a quella Hamiltoniana.

\begin{proposition}{Trasformata di Legendre}{trasformata di Legendre}
    Si consideri uno spaziotempo delle configurazioni $T:\V^{n+1}\to E^1$ per un sistema fisico descritto dalla Lagrangiana $\L:A(\V^{n+1})\to\R$.\\
    A questi dati è associata una mappa liscia $\Leg_\L : A(\V^{n+1})\to F(\V^{n+1})$ detta \bemph{trasformazione di Legendre} definita in coordinate locali come 
    \[ \Leg_\L(t,q,\dot q) = \left( t, q, \frac{\partial \L(t,q\dot q)}{\partial \dot q^1},\ ...\ ,\   \frac{\partial \L(t,q\dot q)}{\partial \dot q^n}\right) \]
    Nel caso in cui la Lagrangiana è della forma
    \[ \L(t,q,\dot q) = \sum_{k,l}A_{k,l}(t,q)\dot q^k \dot q^l + \sum_k \tilde B_k(t,q) \dot q^k + \tilde C(t,q) \]
    Con $A_{k,l}$, $\tilde B_k$ e $\tilde C$ funzioni lisce e la matrice $A = (A_{k,l})_{k,l}$ invertibile, la trasformazione di Legendre è anche un diffeomorfismo.
    \proof 
    Abbiamo già visto che in coordinate locali i momenti coniugati si trasformano come le derivate parziali sotto i cambiamenti di coordinate, dunque le espressioni locali si assemblano in una trasformazione liscia globale.\\
    Sotto le ipotesi date, si ha l'espressione per i $p_k$ data da:
    \[ p^k = \sum_l 2A_{k,l}(t,q)\dot q^l + \tilde B_k(t,q) \]
    Che può essere invertita come 
    \[ \dot q^k(t,q,p) = \sum_{l}\frac{1}{2}(A^{-1})^{k,l}(t,q)(p_l - \tilde B_l(t,q)) \]
    Definendo localmente un'inversa liscia della trasformazione di Legendre.
\end{proposition}

Notiamo che la richiesta sulla forma della Lagrangiana è esattamente quella che avevamo imposto per dimostrare la buona positura delle equazioni di Eulero--Lagrange. Tale condizione in particolare è soddisfatta dalle Lagrangiane di forma standard $\L(t,q,\dot q) = \T_R(t,q,\dot q) - \U(t,q)$, dove $\T_R$ è l'energia cinetica e $\U$ un potenziale.\\
Più in generale, si dice che una Lagrangiana è \bemph{non degenere} se, in ogni scelta di coordinate locali naturali, la famiglia di equazioni 
\[ p_k(t,q,\dot q) = \frac{\partial\L(t,q,\dot q)}{\partial\dot q^k} \]
Può essere invertita ottentendo espressioni $\dot q^k(t,q,p)$ lisce.\\
Ovviamente la trasformata di Legendre è un diffeomorfismo in questo caso più generale di Lagrangiane non degeneri.

\newpage
\section{Equazioni di Hamilton}

In questa parte formuleremo e studieremo le equazioni di Hamilton, che dimostreremo essere equivalenti a quelle di Eulero--Lagrange sotto l'ipotesi di non degenerazione della Lagrangiana.\\
Mentre nella meccanica Lagrangiana la dinamica è codificata dalla scelta di una funzione globale $\L:A(\V^{n+1})\to\R$, nella meccanica Hamiltoniana viene codificata in modo simile dalla scelta di Hamiltoniana $\H$.\\
Il tipo di oggetto che descrive un'Hamiltoniana globale su $F(\V^{n+1})$ però è più sofisticato, in quanto un'Hamiltoniana non è una funzione liscia globale su $F(\V^{n+1})$ a differenza della Lagrangiana. Per comprenderne meglio la costruzione è bene trasformare il dato globale della Lagrangiana $\L$ in una collezione di dati locali.\\
Una Lagrangiana è una mappa globale $\L:A(\V^{n+1})\to\R$, la cui restrizione alle carte adattate $\psi(U)\times\R^n$ definisce una famiglia di mappe lisce locali $\{ \L_U:\psi(U)\times\R^n\to\R \}_{(U,\psi)}$ parametrizzata dalle carte locali adattate a $T:\V^{n+1}\to E^1$.\\
Per costruzione, queste mappe locali devono andare d'accordo sulle intersezioni, ovvero per carte $(U,\psi)$ e $(U,\psi')$ si abbia $\L_{U'}(t',q',\dot q') = \L_U(t,q,\dot q)$ dove $(t,q,\dot q)\mapsto (t',q',\dot q')$ è il cambiamento di coordinate.\\
Viceversa, qualsiasi famiglia di mappe lisce locali $\{\L_U\}_{(U,\psi)}$ che soddisfi queste condizioni definisce globalmente per incollamento una mappa liscia $\L : A(\V^{n+1})\to\R$.

\begin{definition}{Hamiltoniana}{Hamiltoniana}
    Un'Hamiltoniana $\H$ sullo spaziotempo delle fasi $F(\V^{n+1})$ è data da una famiglia di mappe $\{\H_U: \psi(U)\times\R^n\to\R\}_{(U,\psi)}$ tale che per ogni coppia di carte locali $(U,\psi)$ e $(U',\psi')$ con cambiamento di coordinate $(t,q,p)\mapsto (t',q',p')$ si abbia 
    \[ \H_{U'}(t',q',p')  = \H_U(t,q,p) +\sum_k \frac{\partial q'^k}{\partial t}p'_k \]
\end{definition}

\begin{proposition}{Costruzione di un'Hamiltoniana da una Lagrangiana}{Lagrangiana to Hamiltoniana}
    Sia $T:\V^{n+1}\to E^1$ lo spaziotempo delle configurazioni per un sistema con vincoli olonomi e sia $\L:A(\V^{n+1})\to\R$ una Lagrangiana le cui restrizioni a ciascuna carta locale adattata siano tutte non degeneri, ovvero tali che le equazioni $p_k(t,q,\dot q)$ siano localmente invertibili in espressioni $\dot q^k(t,q,p)$.\\
    Allora la famiglia di mappe 
    \[ \left\{\H_U(t,q,p) := \sum_k \dot q^k(t,q,p)p_k - \L_U(t,q,\dot q(t,q,p)) \right\}_{(U,\psi)} \]
    Definisce un'Hamiltoniana su $F(\V^{n+1})$.
    \proof
    Dobbiamo dimostrare che per ogni coppia di carte locali $(U,\psi)$ e $(U',\psi')$ con intersezione non vuota l'Hamiltoniana si trasformi nel modo corretto; ricordando le leggi di trasformazione 
    \[ t' = t \quad q'^k = q'^k(t, q) \quad \dot q'^k = \frac{\partial q'^k(t,q)}{\partial t} + \sum_l \frac{\partial q'^k(t,q)}{\partial q^l}\dot q^l \quad p'_k = \sum_l \frac{\partial q^l}{\partial q'^k}p_k \]
    Segue da un semplice calcolo che 
    \begin{align*}
        \H_{U'}(t',q',p') & = \sum_k\dot q'^k p'_k - \L_{U'}(t',q',\dot q') \\
        & = \sum_k \frac{\partial q'^k}{\partial t}p'_k + \sum_{k,l,m} \frac{\partial q'^k}{\partial q^l}\dot q^l \frac{\partial q^m}{\partial q'^k} p_m - \L_U(t,q,\dot q) \\
        & = \sum_k \frac{\partial q'^k}{\partial t}p'_k + \sum_{l,m}\underbrace{\left( \sum_k \frac{\partial q'^k}{\partial q^l}\frac{\partial q^m}{\partial q'^k} \right)}_{= \frac{\partial q^m}{\partial q^l} = \delta_l^m }\dot q^l p_m - \L_U(t,q,\dot q) \\
        & = \sum_k \frac{\partial q'^k}{\partial t}p'_k + \sum_l \dot q^l p_l - \L_U(t,q,\dot q) = \sum_k \frac{\partial q'^k}{\partial t}p'_k + \H_U(t,q,p)
    \end{align*}
    \qed
\end{proposition}

Notiamo che questa formula per l'Hamiltoniana non è nuova, l'avevamo già vista nel teorema di Jacobi; in casi particolari, $\H$ può essere interpretata come l'energia totale del sistema.\\
La condizione di non degenerazione che abbiamo richiesto in questa proposizione è precisamente la condizione necessaria affinchè la trasformazione di Legendre sia un diffeomorfismo, che è ciò che ci serve per esprimere le Hamiltoniane locali in termini delle coordinate locali di $F(\V^{n+1})$.\\
La bizzarra legge di trasformazione delle Hamiltoniane è imposta dal modo in cui sono costruite rispetto alla Lagrangiana globale $\L$. In un contesto geometrico più avanzato, questa legge di trasformazione viene espressa in termini di una sezione globale di un fibrato affine canonico unidimensionale sullo spaziotempo delle fasi, ma questo va decisamente oltre gli scopi di questo corso.\\
La proposizione precedente ammette un'inversione.

\begin{proposition}{Costruzione di una Lagrangiana da un'Hamiltoniana}{Hamiltoniana to Lagrangiana}
    Sia $T:\V^{n+1}\to E^1$ lo spaziotempo per un sistema con vincoli olonomi e sia $\H = \{\H_U\}_{(U,\psi)}$ un'Hamiltoniana su $F(\V^{n+1})$ tale che tutte le $\H_U(t,q,p)$ siano non degeneri, ovvero le espressioni 
    \[ \dot q^k(t,q,p) = \frac{\partial \H_U(t,q,p)}{\partial p_k} \]
    Siano invertibili in espressioni $p_k = p_k(t,q,\dot q)$.\\
    Allora la famiglia di mappe locali
    \[ \left\{ \L_U(t,q,\dot q) = \sum_k \dot q^k p_k(t,q,\dot q) - \H_U(t,q,p(t,q\dot q)) \right\}_{(U,\psi)} \]
    Definisce una Lagrangiana globale su $A(\V^{n+1})$.
    \proof 
    Dobbiamo dimostrare che $\L_U(t,q,\dot q) = \L_{U'}(t',q',\dot q')$ sull'intersezione non vuota di ogni coppia di carte locali. Con le stesse leggi di trasformazione della proposizione precedente, vediamo 
    \begin{align*}
        \L_U'(t',q',\dot q') & = \sum_k \dot q'^k p'_k - \underbrace{ \sum_k \frac{\partial q'^k}{\partial t}p'_k - \H_U(t,q,p) }_{\H_U'(t',q',p')} \\
        & = \sum_{k} \frac{\partial q'^k}{\partial t}p'_k + \sum_{k,l,m} \frac{\partial q'^k}{\partial q^l}\dot q^l \frac{\partial q^m}{\partial q'^k} p_m - \H_U(t,q,p) - \sum_k \frac{\partial q'^k}{\partial t}p'_k \\
        & = \sum_{l,m} \underbrace{\left( \sum_k \frac{\partial q'^k}{\partial q^l}\frac{\partial q^m}{\partial q'^k} \right)}_{= \frac{\partial q^m}{\partial q^l} = \delta_l^m }\dot q^l p_m - \H_U(t,q,p) \\
        & = \sum_l \dot q^l p_l - \H_U(t,q,p) = \L_U(t,q,p)
    \end{align*}
    \qed
\end{proposition}

Come vorremmo sperare, abbiamo il seguente teorema.

\begin{theorem}{Equivalenza tra Lagrangiana e Hamiltoniana}{}
    Le due costruzioni precedenti sono inverse l'una dell'altra e stabiliscono una biezione tra le Lagrangiane non degeneri su $A(\V^{n+1})$ e le Hamiltoniane non degeneri su $F(\V^{n+1})$.
    \proof 
    Dobbiamo dimostrare due cose: $(1)$ che le due costruzioni preservano la non-degenerazione, e $(2)$ che sono inverse l'una all'altra.
    \begin{enumerate}
        \item \begin{itemize}
            \item Osserviamo che nel caso $\L\mapsto \H$ si ha 
                \begin{align*}
                    \frac{\partial\H_U(t,q,p)}{\partial p_k} & = \frac{\partial}{\partial p_k}\left( \sum_l \dot q^l (t,q,p) p_l - \L_U(t,q,\dot q(t,q,p)) \right) \\
                    & = \sum_l \frac{\partial \dot q^l}{\partial p_k}p_l + \dot q^k(t,q,p) +\sum_l \underbrace{\frac{\partial \L_U(t,q,\dot q(t,p,q))}{\partial \dot q^l}}_{=: p_l} \frac{\partial \dot q^l}{\partial p_k} = \dot q^k(t,q,p)
                \end{align*}
                Che è invertibile per 
                \[ p_k(t,q,\dot q) = \frac{\partial \L_U(t,q, \dot q)}{\partial \dot q^k} \]
            \item Nel caso $\H\mapsto \L$ invece abbiamo 
                \begin{align*}
                    \frac{\partial\L_U(t,q,\dot q)}{\partial \dot q^k} & = \frac{\partial }{\partial \dot q^k}\left( \sum_l \dot q^l p_l - \H_U(t,q,p(t,q,\dot q)) \right) \\
                    & = p_k(t,q,\dot q)+ \sum_l \dot q^l\frac{\partial p_l(t,q,\dot q)}{\partial \dot q^k} - \sum_l \underbrace{\frac{\partial \H_U(t,q,p(t,q,\dot q))}{ \partial p_l}}_{=: \dot q^l} \frac{\partial p_l(t,q,\dot q)}{\partial \dot q^k} = p_k(t,q,\dot q)
                \end{align*}
                Che è invertibile per 
                \[ \dot q^k(t,q,p) = \frac{\partial \H_U(t,q,p)}{\partial p_k} \]
        \end{itemize}
        \item \begin{itemize}
            \item Data una Lagrangiana non degenere $\{\L_U\}$ troviamo la Lagrangiana associata alla sua Hamiltoniana $\{\H_U\}$:
                \[ \sum_k \dot q^k p_k -\H_U = \sum_k \dot q^k p_k - \sum_k \dot q^k p_k + \L_U = \L_U \]
            \item Analogamente per un'Hamiltoniana non degenere $\{H_U\}$: 
                \[ \sum_k \dot q^k p_k -\L_U = \sum_k \dot q^kp_k - \sum_k \dot q^k p_k + \H_U = \H_U \]
        \end{itemize}
    \end{enumerate}
    \qed 
\end{theorem}

Possiamo ora descrivere come si presentano le equazioni di Eulero--Lagrange quando formulate dal punto di vista dell'Hamiltoniana invece che della Lagrangiana; per farlo sarà comodo esprimerle in una forma equivalente.\\
Consideriamo una generica linea di universo a valori in una carta locale $\Gamma_A:I\to\psi(U)\times \R^n \subset A(\V^{n+1})$ che mappa $t\mapsto (t,q(t),\dot q(t))$. Questa sorge come estensione canonica $\Gamma_A = \hat\Gamma$ di una linea di universo $\Gamma:I\to\psi(U)$ se e solo se per ogni $k$ vale
\[ \dot q^k(t) = \frac{\di q^k(t)}{\di t} \]
Allora possiamo riscrivere le equazioni di Eulero--Lagrange per nel sistema di equazioni 
\[ \begin{dcases}
    \frac{\di }{\di t}\frac{\partial \L(t,q(t),\dot q(t))}{\partial \dot q^k} - \frac{\partial \L_U(t, q(t),\dot q(t))}{\partial q^k} \\ \dot q^k(t) = \frac{\di q^k(t)}{\di t}
\end{dcases} \]

\begin{theorem}{Equazioni di Hamilton}{equazioni di Hamilton}
    Consideriamo una Lagrangiana non degenere $\{\L_U\}$ e una linea di universo locale $\Gamma_A : I \to \psi(U)\times \R^n\subset A(\V^{n+1})$ che mappa $t\mapsto (t,q(t),\dot q(t))$.\\
    Usando la trasformazione di Legendre, questi dati ammettono una descrizione equivalente in termini di un'Hamiltoniana $\{\H_U\}$ e una linea di universo locale $\Gamma_F:I \to\psi(U)\times \R^n \subset F(\V^{n+1})$ che mappa $t\mapsto (t,q(t),p(t))$.\\
    Allora $\Gamma_A$ soddisfa le equazioni di Eulero--Lagrange nella forma equivalente di cui sopra se e solo se $\Gamma_F$ soddisfa le \bemph{equazioni di Hamilton}:
    \[ \begin{dcases}
        \frac{\di p_k(t)}{\di t} = -\frac{\partial \H_U(t,q(t),p(t))}{\partial q^k}\\
        \frac{\di q^k(t)}{\di t} = \frac{\partial \H_U(t,q(t),p(t))}{\partial p_k}
    \end{dcases} \]
    \proof 
    Dato che la trasformata di Legendre è un diffeomorfismo nel caso non degenere, basta provare che le equazioni di Hamilton implicano quelle di Eulero--Lagrange: vediamo che per la definizione di Hamiltoniana 
    \[ \H_U(t,q,p) = \sum_k \dot q^k(t,q,p) p_k -\L_U(t,q,\dot q(t,q,p)) \]
    Si ha per il lato destro
    \begin{align*}
        -\frac{\partial \H_U(t,q,p)}{\partial q^k} & = -\sum_l \frac{\partial \dot q^l(t,q,p)}{\partial q^k}p_l +\frac{\partial \L_U(t,q,\dot q(t,q,p))}{\partial q^k} + \sum_l \underbrace{\frac{\partial \L_U(t,q,\dot q(t,q,p))}{\partial \dot q^l}}_{=: p_l}\frac{\partial \dot q^l}{\partial q^k}\\ 
        & = \frac{\partial \L_U(t,q,\dot q(t,q,p))}{\partial q^k} \\
        \frac{\partial\H_U(t,q,p)}{\partial p_k} & = \sum_l \frac{\partial \dot q^l(t,q,p)}{\partial p^k}p_l + \dot q^k(t,q,p) -\sum_l \underbrace{\frac{\partial \L_U(t,q,\dot q(t,q,p))}{\partial \dot q^l}}_{=: p_l}\frac{\partial \dot q^l(t,q,\dot q(t,q,p))}{\partial p_k} \\
        & = \dot q^k(t,q,p)
    \end{align*}
    Valutando sulla trasformata di Legendre di $\Gamma_A(t) = (t,q(t),\dot q(t))$ data da 
    \[ \Gamma_F(t):= (t,q,p(t)) = \left( t, q^1(t),\ ...\ ,\ q^n(t),\frac{\partial \L_U(t,q(t)\dot q(t))}{\partial \dot q^1},\ ...\ ,\  \frac{\partial \L_U(t,q(t),\dot q(t))}{\partial\dot q^n} \right) \]
    Otteniamo le equazioni di Eulero--Lagrange: 
    \begin{align*}
        \frac{\di}{\di t}\frac{\partial \L_U(t,q(t),\dot q(t))}{\partial \dot q^k} & = \frac{\di p_k(\Gamma_F(t))}{\di t} = - \frac{\partial \H_U(\Gamma_F(t))}{\partial q_k} \\
        & = \frac{\partial \L(t,q(t),\dot q(t))}{\partial q^k} \\
        \frac{\di q^k(t)}{\di t} &  = \frac{\partial \H_U(\Gamma_F(t)) }{\partial p_k} = \dot q^k(t)
    \end{align*}
    \qed
\end{theorem}

Sebbene questo teorema sia enunciato solo nel caso locale, come abbiamo visto qualche pagina fa possiamo estenderlo al caso globale: data una linea di universo $\Gamma_A:I\to A(\V^{n+1})$ con trasformata di Legendre $\Gamma_F:I\to F(\V^{n+1})$, possiamo ricoprire $I$ in sottointervalli $\{I_U\}$ tali che le restrizioni delle linee di universo si fattorizzino in carte locali:
\[\begin{tikzcd}
	{I_U} && {A(\V^{n+1})} && {I_U} && {F(\V^{n+1})} \\
	\\
	&& {\psi(U)\times\R^n} &&&& {\psi(U)\times \R^n}
	\arrow["{\Gamma_A|_{I_U}}"{description}, from=1-1, to=1-3]
	\arrow["{\Gamma_A}"{description}, dashed, from=1-1, to=3-3]
	\arrow["{\Gamma_F|_{I_U}}"{description}, from=1-5, to=1-7]
	\arrow["{\Gamma_F}"{description}, dashed, from=1-5, to=3-7]
	\arrow[hook, from=3-3, to=1-3]
	\arrow[hook, from=3-7, to=1-7]
\end{tikzcd}\]
E applicando il teorema a tutta questa famiglia di linee di universo locali otteniamo l'equivalenza globale.\\
Notiamo che le equazioni di Hamilton costituiscono un sistema di equazioni del primo ordine in forma normale per $q^k(t)$ e $p_k(t)$. Questo sistema tuttavia le tratta in modo molto più simmetrico di quanto la meccanica Lagrangiana tratti $q(t)$ e $\dot q(t)$; in particolare vengono preservate dalla trasformazione $(q^k,p_k)\mapsto (-p_k, q_k)$, che è un esempio di una cosiddetta \bemph{trasformazione canonica}, argomento che non tratteremo in questo corso ma che è molto importante per gli sviluppi della meccanica Hamiltoniana\footnote{In effetti, questa equivalenza tra \textit{posizione} e \textit{momento} di una \textit{particella} potrebbe far venire qualcosa in mente ad alcuni lettori}.\\

\newpage
\section{Parentesi di Poisson}

In questa parte introduciamo la cosiddetta parentesi di Poisson, una struttura algebrica definita canonicamente su $F(\V^{n+1})$. Questa non serve solo a riscrivere in una forma molto coincisa le equazioni di Hamilton, ma anche a formulare il teorema di Noether per la meccanica Hamiltoniana e affrontare il problema della quantizzazione in meccanica quantistica.

\begin{definition}{Parentesi di Poisson locale}{parentesi di Poisson locale}
    Sia $(U,\psi)$ una carta locale adattata alla sommersione suriettiva $T:\V^{n+1}\to E^1$ e si consideri la carta locale associata $\psi(U)\times\R^n$ per $F(\V^{n+1})$.\\
    La \bemph{parentesi di Poisson} di due mappe lisce $f_U, g_U : \psi(U)\times\R^n \to \R $ è data da 
    \[ \langle f_U, g_U\rangle_U = \sum_k \frac{ \partial f_U }{\partial q^k} \frac{\partial g_U}{\partial p_k} - \frac{\partial f_U}{\partial p_k}\frac{\partial g_U}{\partial q^k} \]
\end{definition}

\begin{proposition}{parentesi di Poisson globale}{parentesi di Poisson globale}
    Le parentesi di Poisson locali $\langle\cdot, \cdot\rangle_U$ si assemblano in una parentesi di Poisson globale $\langle\cdot,\cdot\rangle:\C^\infty (F(\V^{n+1}))\times \C^\infty (F(\V^{n+1})) \to \C^\infty (F(\V^{n+1}))$ sull'algebra delle funzioni reali lisce sullo spaziotempo delle fasi.\\
    In particolare, ricordando che ogni mappa liscia $h:F(\V^{n+1})\to \R$ ammette una descrizione equivalente in termini di una famiglia di mappe locali lisce $\{h_U\}_{(U,\psi)}$ che concordano sull'intersezione delle carte, la parentesi di Poisson globale $\langle f,g\rangle$ è definita dalla famiglia $\{ \langle f_U, g_U\rangle_U \}_{(U,\psi)}$
    \proof 
    Dobbiamo dimostrare che sull'intersezione $U\cap U'$ le parentesi locali soddisfano $\langle f_{U'}, g_{U'} \rangle_{U'}(t',q',p') = \langle f_U, g_U\rangle_U(t,q,p)$, ricordando la legge di trasformazione 
    \[ t' = t, \quad q' = q'(t,q), \quad p'_k = \sum_l \frac{\partial q^l}{\partial q'^k}p_l \]
    E usando il fatto che $f_U$ e $g_U$ concordano con $f_{U'}$ e $g_{U'}$ sull'intersezione otteniamo la regola della catena:
    \begin{align*}
        \langle f_U, g_U\rangle_U(t,q,p) & = \sum_k \frac{ \partial f_U }{\partial q^k} \frac{\partial g_U}{\partial p_k} - \frac{\partial f_U}{\partial p_k}\frac{\partial g_U}{\partial q^k} \\
        & = \sum_{k,l,m} \left[ \left( \frac{\partial q'^l}{\partial q'^k} \frac{\partial f_{U'}}{\partial q'^l} + \frac{\partial p'_l}{\partial q^k}\frac{\partial f_{U'}}{\partial p'_l} \right)\underbrace{\frac{\partial p'_m}{\partial p_k}}_{= \frac{\partial q^k}{\partial q'^m}}\frac{\partial g_{U'} }{\partial p'_m} - \underbrace{\frac{\partial p'_l}{\partial p_k}}_{=\frac{\partial q^k}{\partial q'^l}}\frac{\partial f_{U'}}{\partial p'_l}\left( \frac{ \partial q'^m }{\partial q^k} \frac{\partial g_{U'}}{\partial q'^m} + \frac{\partial p'_m}{q^k}\frac{\partial g_{U'}}{\partial p'_m} \right) \right] \\
        & = \sum_k \left(\frac{ \partial f_{U'} }{\partial q'^k} \frac{\partial g_{U'}}{\partial p'_k} - \frac{\partial f_{U'}}{\partial p'_k}\frac{\partial g_{U'}}{\partial q'^k}\right) + \underbrace{\sum_{k,l,m} \left( \frac{\partial p'_l}{\partial q^k}\frac{\partial q^k}{\partial q'^m}- \frac{\partial p'_m}{\partial q^k}\frac{\partial q^k}{\partial q'^l} \right)\frac{\partial f_U}{\partial p'_l}\frac{\partial g_U}{\partial p'_m}}_{= 0} \\
        & = \langle f_{U'}, g_{U'} \rangle_{U'}(t',q',p')
    \end{align*}
    Dove nell'ultimo passaggio abbiamo usato l'identità
    \[ \sum_k \frac{\partial q'^l}{\partial q^k}\frac{\partial q^k}{\partial q'^m} = \frac{\partial q'^l}{\partial q'^m} = \delta_m^l \]
    E il termine che si annulla viene dal fatto che
    \[\sum_k \frac{\partial p'_l}{\partial q^k}\frac{\partial q^k}{\partial q'^m} = \sum_{k,s} p_s \frac{\partial}{\partial q^k}\left( \frac{\partial q^s}{\partial q'^l} \right)\frac{\partial q^k}{\partial q'^m} = \sum_{k,s,u} p_s \frac{\partial^2 q^s}{\partial q'^u \partial q'^l}\frac{ \partial q'^u}{\partial q^k}\frac{\partial q^k}{\partial q'^m} = \sum_s p_s \frac{ \partial^2 q^s }{\partial q'^l \partial q'^m}\]
    È simmetrico rispetto allo scambio di $l$ e $m$.
    \qed
\end{proposition}

Notiamo che se introduciamo la matrice 
\[ S := \begin{pmatrix}
    \mathbf{0}_{n\times n} & \id_{n\times n} \\ - \id_{n\times n} & \mathbf{0}_{n\times n}
\end{pmatrix} \]
La parentesi di Poisson locale può essere scritta più compattamente come 
\[ \langle f_U, g_U \rangle_U = \sum_{a,b} S^{a,b} \frac{\partial f_U}{\partial z^a}\frac{\partial g_U}{\partial z^b} \]
Forma occasionalmente utile per i calcoli.\\
Le parentesi di Poisson delle coordinate locali assumono i valori 
\[ \langle q^k, q^l \rangle_U = 0,\quad \langle p_k, p_l\rangle_U = 0 \quad \langle q^k, p_l\rangle_U = \delta_l^k  \]
Queste identità sono un possibile punto di partenza per passare dalla meccanica classica a quella quantistica. Infatti sostituendo $q$ e $p$ con operatori, ad esempio su uno spazio di Hilbert complesso, $\hat q$ e $\hat p$ i cui commutatori soddisfano le identità analoghe
\[ [\hat q^k, \hat q^l]= 0,\quad [\hat p_k, \hat p_l] = 0,\quad [\hat q^k,\hat p_l] = i\hbar \delta_l^k  \]
Dove $\hbar$ è la costante di Planck ridotta. La matematica che sorge da questa transizione è estremamente interessante, ma purtroppo non abbiamo il tempo di esplorarla in questo corso.

\begin{proposition}{Proprietà della parentesi di Poisson}{proprietà della parentesi di Poisson}
    La parentesi di Poisson gode delle seguenti proprietà:\begin{enumerate}
        \item \bemph{bilinearità}, ovvero 
        \[ \langle \alpha f + \beta g, h \rangle = \alpha\langle f,h \rangle+\beta \langle g, h\rangle,\qquad \langle f, \alpha g + \beta h \rangle = \alpha\langle f,g\rangle + \beta\langle f, h\rangle \]
        \item \bemph{Antisimmetria}, ovvero $\langle f,g\rangle = - \langle g, f\rangle$.
        \item \bemph{Identità di Jacobi}\footnote{Legata alle permutazioni cicliche $fgh\to hfg \to ghf$}, ovvero 
        \[ \langle\langle f, g \rangle, h\rangle + \langle\langle g,h\rangle,f\rangle + \langle\langle h,f\rangle, g\rangle = 0 \]
        \item \bemph{Biderivazione}, o regola di Leibniz:
        \[ \langle fg, h \rangle = f \langle g,h\rangle + \langle f,h\rangle g, \qquad \langle f, gh\rangle = \langle f,g\rangle h + g\langle f,h\rangle \]
    \end{enumerate}
    \proof 
    È sufficiente dimostrare localmente queste proprietà, dunque dall'espressione locale vista sopra 
    \[ \langle f_U, g_U\rangle_U = \sum_{a,b}S^{a,b} \frac{\partial f_U}{\partial z^a} \frac{\partial g_U}{\partial z^b}  \]
    Seguono immediatamente $(1)$, $(2)$ e $(4)$ dalle proprietà della derivata e dall'antisimmetria di $S$.\\
    Rimane da mostrare la $(3)$: esplicitando il primo termine e omettendo $U$ per alleggerire la notazione otteniamo:
    \begin{align*}
        \langle\langle f, g \rangle, h\rangle & = \sum_{a,b}S^{a,b}\frac{\partial \langle f, g \rangle}{\partial z^a}\frac{\partial h}{\partial z^b} = \sum_{a,b,c,d}S^{a,b}S^{c,d}\frac{\partial }{\partial z^a}\left( \frac{\partial f}{\partial z^c}\frac{\partial g}{\partial z^d} \right)\frac{\partial h}{\partial z^b} \\ 
        & = \sum_{a,b,c,d} S^{a,b}S^{c,d}\left( \frac{\partial^2 f}{\partial z^a \partial z^c} \frac{\partial g}{\partial z^d} \frac{\partial h}{\partial z^b} + \frac{\partial f}{\partial z^c}\frac{\partial^2 g}{\partial z^a\partial z^d}\frac{\partial h}{\partial z^b} \right) 
    \end{align*}
    Quindi possiamo riscrivere il lato sinistro dell'identità di Jacobi come 
    \begin{align*}
        \sum_{a,b,c,d}S^{a,b}S^{c,d} & \left( \underbrace{\frac{\partial^2 f}{\partial z^a \partial z^c} \frac{\partial g}{\partial z^d} \frac{\partial h}{\partial z^b}}_{=: \heartsuit_1} + \underbrace{\frac{\partial f}{\partial z^c}\frac{\partial^2 g}{\partial z^a\partial z^d}\frac{\partial h}{\partial z^b}}_{=:\diamondsuit_1} \right. \\
        & + \underbrace{\frac{\partial^2 g}{\partial z^a \partial z^c} \frac{\partial h}{\partial z^d} \frac{\partial f}{\partial z^b}}_{=: \diamondsuit_2} + \underbrace{\frac{\partial g}{\partial z^c}\frac{\partial^2 h}{\partial z^a\partial z^d}\frac{\partial f}{\partial z^b}}_{=:\clubsuit_1} \\
        & \left. + \underbrace{\frac{\partial^2 h}{\partial z^a \partial z^c} \frac{\partial f}{\partial z^d} \frac{\partial g}{\partial z^b}}_{=: \clubsuit_2} + \underbrace{\frac{\partial h}{\partial z^c}\frac{\partial^2 f}{\partial z^a\partial z^d}\frac{\partial g}{\partial z^b}}_{=:\heartsuit_2} \right)
    \end{align*} 
    Si vede con un poco di lavoro che i termini con lo stesso seme\footnote{a dire il vero, una volta messi in sommatoria} sono opposti l'uno all'altro: consideriamo $\heartsuit$, per gli altri semi il ragionamento è analogo:
    \begin{align*}
        \heartsuit_2 & = \sum_{a,b,c,d} S^{a,b}S^{c,d} \frac{\partial h}{\partial z^c}\frac{\partial^2 f}{\partial z^a\partial z^d}\frac{\partial g}{\partial z^b} \\
        & = \sum_{a,b,c,d} S^{c,d}\underbrace{S^{a,b}}_{= -S^{b,a}} \frac{\partial h}{\partial z^a}\frac{\partial^2 f}{\underbrace{\partial z^c\partial z^b}_{= \partial z^b \partial z^c}}\frac{\partial g}{\partial z^d} \\
        & = -\sum_{a,b,c,d} S^{b,a}S^{c,d} \frac{\partial^2 f}{\partial z^b \partial z^c } \frac{\partial g}{\partial z^d}\frac{\partial h}{\partial z^a} \\
        & = -\sum_{a,b,c,d} S^{a,b}S^{c,d} \frac{\partial^2 f}{\partial z^a \partial z^c } \frac{\partial g}{\partial z^d}\frac{\partial h}{\partial z^b} = - \heartsuit_1
    \end{align*}
    Dove nel secondo passaggio abbiamo rinominato $a,b$ in $c,d$ e viceversa e nel quarto rinominato $a$ in $b$ e viceversa.
    \qed
\end{proposition}

Le proprietà $(1)$, $(2)$ e $(3)$ sono le proprietà che caratterizzano una cosiddetta \bemph{algebra di Lie}; pertanto la parentesi di Poisson dota $\C^{\infty}(F(\V^{n+1}))$ di una struttura di algebra di Lie con una proprietà aggiuntiva, la $(4)$, che la rende compatibile con la moltiplicazione puntuale di funzioni. Nel contesto della geometria differenziale, una varietà dotata di una parentesi di Poisson è detta \bemph{varietà di Poisson}.

\begin{proposition}{Equazioni di Hamilton in forma di parentesi di Poisson}{parentesi di poisson per le equazioni di Hamilton}
    Nel contesto della definizione di cui sopra, le equazioni di Hamilton possono essere riscritte equivalentemente come 
    \[ \begin{dcases}
        \frac{\di p_k(t)}{\di t} = \langle p_k, \H_U\rangle_U(t,q(t),p(t)) \\
        \frac{\di q^k(t)}{\di t} = \langle q^k, \H_U\rangle_U(t,q(t),p(t))
    \end{dcases} \]
    \proof 
    Segue per calcolo diretto, infatti per $p_k$ abbiamo 
    \[ \langle p_k, \H_U\rangle_U = \sum_l \left( \underbrace{\frac{\partial p_k}{\partial q^l}}_{=0}\frac{\partial \H_U}{\partial p_l} - \underbrace{\frac{\partial p_k}{\partial p_l}}_{=\delta_l^k}\frac{\partial \H_U}{\partial q^l} \right) = -\frac{\partial \H_U}{\partial q^k} \]
    E analogamente per $q^k$
    \qed 
\end{proposition}

La parentesi di Poisson non è solo utile per riscrivere in maniera più compatta le equazioni di Hamilton, ma anche per studiare gli integrali primi del moto.

\begin{definition}{Integrale primo del moto}{integrale primo}
    Si consideri un sistema fisico descritto da uno spaziotempo delle configurazioni $T:\V^{n+1}\to E^1$ e un'Hamiltoniana $\H = \{\H_U\}_{(U,\psi)}$ sullo spaziotempo delle fasi $F(\V^{n+1})$.\\
    Un \bemph{integrale primo del moto} è una funzione liscia globale $f:F(\V^{n+1})\to \R$ tale che in ogni carta locale $\psi(U)\times\R^n$ e lungo ogni soluzione $(t,q(t),p(t))$ locale delle equazioni di Hamilton si abbia 
    \[\frac{\di f_U(t,q(t),p(t))}{\di t} = 0\]
\end{definition}

\begin{proposition}{Caratterizzazione degli integrali primi del moto}{caratterizzazione degli integrali primi del moto}
    Nel contesto della definizione di cui sopra, $f = \{ f_U\}_{(U,\psi)}$ è un integrale primo del moto se e solo se in ogni carta locale $\psi(U)\times\R^n$ si ha 
    \[ \frac{\partial f_U}{\partial t} + \langle f_U,\H_U\rangle_U = 0 \]
    \proof 
    Usando le equazioni di Hamilton e sottintendendo che tutte le funzioni sono da valutare lungo $(t,q(t),p(t))$ in $U$ possiamo scrivere 
    \begin{align*}
        \frac{\di f}{\di t} & = \frac{\partial f}{\partial t} + \sum_k\left( \frac{\partial f}{\partial q^k} \frac{\di q^k}{\di t} + \frac{\partial f}{\partial p_k}\frac{\di p_k}{\di t} \right) \\
        & = \frac{\partial f}{\partial t} + \sum_k\left( \frac{\partial f}{\partial q^k} \frac{\partial \H}{\partial p_k} - \frac{\partial f}{\partial p_k}\frac{\partial \H}{\partial q^k} \right) \\
        & = \frac{\partial f}{\partial t} + \langle f, \H\rangle
    \end{align*}
    Poichè le equazioni di Hamilton sono in forma normale, ogni punto di $F(\V^{n+1})$ funge da dato iniziale, dunque ogni punto della carta appartiene a qualche linea di universo e perciò questa identità è vera su tutta la carta.
\end{proposition}

Dalla sua derivazione, è chiaro che questa condizione equivalente per un integrale primo è interpretabile come l'affermazione che $f$ è invariante rispetto alla dinamica generata dall'Hamiltoniana $\H$ del sistema. Abbiamo già visto sopra che la dinamica delle coordinate del sistema viene codificata dall'azione di $\langle \cdot, \H\rangle$ sulle coordinate, dunque con un po' di lavoro si può mostrare che la dinamica di una funzione liscia $f$ è codificata dall'azione del campo vettoriale, o derivazione 
\[ Z = \frac{\partial}{\partial t} + \langle\cdot, \H\rangle \]
E che dunque gli integrali primi corrispondono agli zeri di questo campo vettoriale, ovvero le funzioni che non cambiano lungo la sua dinamica.\\
Allo stesso modo, si può interpetare questa condizione dicendo che il sistema Hamiltoniano possiede una simmetria generata da $f$ tramite la derivazione $X_f = \langle f,\cdot\rangle$, che può essere integrata a un gruppo a un parametro di automorfismi di $F(\V^{n+1})$.\\
La relazione tra questi due punti di vista viene codificata precisamente nel \bemph{teorema di Noether Hamiltoniano}, che può essere reperito nel capitolo 13.3.3 di \cite{Moretti2024}, ma che purtroppo non abbiamo tempo di trattare.\\
Concludiamo con un bel risultato sulle strutture algebriche definibili sull'insieme degli integrali primi del moto.

\begin{theorem}{Struttura di algebra di Poisson per gli integrali primi}{algebra di Poisson degli integrali primi}
    Indichiamo con 
    \[ \I_\H := \left\{ f \in \C^{\infty}(F(\V^{n+1})) : \frac{\partial f}{\partial t} + \langle f,\H \rangle = 0 \right\} \]
    L'insieme degli integrali primi del moto di un sistema fisico descritto da un'Hamiltoniana $\H$.\\
    Allora $(\I_\H, \langle\cdot,\cdot\rangle )$ è una sottoalgebra di Poisson di $(\C^{\infty}(F(\V^{n+1})), \langle\cdot,\cdot\rangle )$, vale a dire che include la funzione identicamente $1$ ed è chiuso per combinazioni lineari, prodotti e parentesi di Poisson dei suoi elementi.
    \proof 
    È sufficiente provare la proposizione localmente:\begin{itemize}
        \item La chiusura per combinazioni lineari segue dalla linearità della derivata rispetto a $t$ e di $\langle \cdot , \H\rangle$.
        \item Segue dal fatto che la funzione costante $1$ è costante e dalla proprietà di derivazione di $\langle \cdot , \H\rangle$.
        \item Segue dall'antisimmetria e dall'identità di Jacobi.
    \end{itemize}
    \qed
\end{theorem}

Questo risultato può esserci utile per cercare nuovi integrali primi del moto quando ne abbiamo già alcuni: non solo considerando banalmente loro combinazioni lineari o prodotti, ma anche considerando le loro parentesi di Poisson.

\nocite{*}
\printbibliography

\end{document}